%% Paper II of the Gaussian Schmidt arrangement project.
%% Sources: the repository documents phase-kronecker-limit.md and
%% schmidt-units.md (primary), euclidean-moduli-invariants.md SS5.5-5.6 and
%% moduli-invariants.md SS5.4-5.7 (background), and Paper I
%% (papers/1-schmidt-circles/schmidt-circles.tex), cited as [SchmidtI].
%% Machine verification: see Appendix A and scripts/phase_klf.py,
%% scripts/schmidt_units.py.
\documentclass[11pt]{amsart}

\usepackage{amsmath,amssymb,amsthm}
\usepackage{booktabs}
\usepackage{array}
\usepackage[hidelinks]{hyperref}
\usepackage[expansion=false]{microtype}

\allowdisplaybreaks

\numberwithin{equation}{section}

\theoremstyle{plain}
\newtheorem{theorem}{Theorem}[section]
\newtheorem{proposition}[theorem]{Proposition}
\newtheorem{lemma}[theorem]{Lemma}
\newtheorem{corollary}[theorem]{Corollary}
\newtheorem{conjecture}[theorem]{Conjecture}

\theoremstyle{definition}
\newtheorem{definition}[theorem]{Definition}
\newtheorem{example}[theorem]{Example}
\newtheorem{openproblem}[theorem]{Open problem}

\theoremstyle{remark}
\newtheorem{remark}[theorem]{Remark}

%% notation (as in Paper I)
\newcommand{\C}{\mathbb{C}}
\newcommand{\R}{\mathbb{R}}
\newcommand{\Q}{\mathbb{Q}}
\newcommand{\Z}{\mathbb{Z}}
\newcommand{\HH}{\mathbb{H}}
\newcommand{\Zi}{\Z[i]}
\newcommand{\Rhat}{\widehat{\R}}
\newcommand{\SL}{\mathrm{SL}}
\newcommand{\PSL}{\mathrm{PSL}}
\newcommand{\Gal}{\mathrm{Gal}}
\newcommand{\Cl}{\mathrm{Cl}}
\newcommand{\im}{\operatorname{Im}}
\newcommand{\re}{\operatorname{Re}}
\newcommand{\Sarr}{\mathcal{S}}
\newcommand{\ideala}{\mathfrak{a}}
\newcommand{\idealb}{\mathfrak{b}}
\newcommand{\idealc}{\mathfrak{c}}
\newcommand{\idealp}{\mathfrak{p}}
\newcommand{\idealP}{\mathfrak{P}}
\newcommand{\idealq}{\mathfrak{q}}
\newcommand{\idealr}{\mathfrak{r}}
\newcommand{\ideals}{\mathfrak{s}}
\newcommand{\idealm}{\mathfrak{m}}
\newcommand{\GZ}{\mathrm{GZ}}
\newcommand{\Ne}{N_{e}}
\newcommand{\covol}{\operatorname{covol}}
\newcommand{\Lie}{\operatorname{Lie}}
\newcommand{\lat}[1]{\left[#1\right]}
\newcommand{\Lprim}{L_{\mathrm{prim}}}
\newcommand{\Oc}{\mathcal{O}}
\newcommand{\Sig}{\Sigma}
\newcommand{\eps}{\varepsilon}
\newcommand{\PaperI}{\cite{SchmidtI}}

\begin{document}

\title[Schmidt circles and elliptic units]{Schmidt circles and elliptic
units:\\ Kronecker limit formulas and the Robert index}

%% TODO before submission: complete author name and affiliation.
\author{Karl Ivar}
\email{karlivar2004@gmail.com}

\begin{abstract}
The phases $u_f$ of the Gaussian Schmidt arrangement --- the sixth
coordinate on $\SL_2(\Z)\backslash\SL_2(\C)/\SL_2(\Z)$, quantized on the
Bianchi group in the companion paper --- are, classwise, twisted
$\Delta$-quotients at CM points.  We pursue this to its quantitative end.
First, the character sums $S(\chi) = \sum\chi\log|u|$ over the classes of
a level obey Kronecker limit formulas in both aspects of the theory:
$S(\chi) = -2L'(0,\chi) + \tfrac23\Sigma_0(\chi) + \tfrac12\Sigma_{1728}(\chi)$
for the Euclidean phases of the curvature-$2n$ disks (over the ring class
group of $\Z + n\Z[i]$), and, for the hyperbolic phases of discriminant
$1-n^2$, $S(\chi) = 0$ or $-4L'(0,\chi) + \tfrac43\Sigma_0 + \Sigma_{1728}$
according to the value of $\chi$ on the twist class, the unit
$n + \sqrt{n^2-1}$ cancelling identically.  Here $L(s,\chi)$ is the Epstein
$L$-function of the level and $\Sigma_x(\chi) = \sum\chi\log|j - x|$ a
Gross--Zagier collision quantity.  For genus characters everything is in
closed form through Sturm-proved factorizations of $L(s,\chi)$; the real
quadratic field carrying the Euclidean genus sum at conductor $n$ is
$\Q(\sqrt n)$, and the $j$-dressings factor exactly in it with split-prime
support at Gross--Zagier primes.  Second, the $\idealr$-twisted
$\Delta$-ratios $R_f = r_0^{6}\Delta(\idealb_1)/\Delta(\idealr^{-1}\idealb_1)$
are algebraic units for every class of every odd level, imprimitive strata
included --- a local rigidity argument for isogenies attached to invertible
ambiguous ideals, needing no quasi-canonical-lifting input --- and the
Euclidean $\Delta$-data $n^{12}\Delta(\Lambda_\idealc)/\Delta(\Z[i])$ have a
$\idealP$-adic valuation independent of the class, with an explicit value
that settles the split-prime denominator ladders.  A canonical sixth root
$w_f$ obeys the laws at first power; at level $7$ it is $-(2\pm\sqrt3)$.
Third, on the cubic layer of the Euclidean theory the $\Delta$-coset unit
$\theta_u$ has index $[\Oc_{L_3}^\times : \langle -1, \theta_u\rangle] =
8\,h_{L_3}\,C_n(0)$ in the unit group of the real cubic subfield $L_3$ of the
ring class field, certified at all eight levels $n \le 27$ with a cubic
character, with $h_{L_3} = 1, 1, 3, 2$ at the primitive levels $9, 11, 13,
23$, the last out of sample: the class number of the cubic field is the
Robert index of the Schmidt units.  Every finite claim is machine-verified;
the scripts and precision policies are catalogued in an appendix.
\end{abstract}

\maketitle
\tableofcontents

%%% ============================================================
\section{Introduction}\label{sec:intro}

Let $\Sarr = \PSL_2(\Zi)\cdot\Rhat$ be the Gaussian Schmidt arrangement,
the orbit of the extended real line under the Bianchi group
\cite{Schmidt1975, StangeTAMS, StangeIMRN}.  The companion paper
\PaperI{} --- cited throughout as [SchmidtI], with its theorem numbers ---
stratified $\Sarr$ by an integer level, counted it in three geometries,
computed the involution $\sigma(X) = \bar X^{-1}$ on the form classes of a
level, and constructed a \emph{phase}: a sixth coordinate on the double
cosets $\SL_2(\Z)\backslash\SL_2(\C)/\SL_2(\Z)$ which, on Schmidt matrices,
quantizes to algebraic numbers $u_f$ indexed by the primitive form classes
$f$ of discriminant $1 - n^2$, and to Euclidean phases $u_\idealc$ indexed
by the ring class group $\Cl(\Z + n\Zi)$ of the curvature-$2n$ disks.  Its
last structural statement was the closed form
\[
  u_f = -\,\eps\,\mu^{-2}\,\frac{h_2(\idealb_1)}{h_2(\idealb_2)},
  \qquad
  u_\idealc^{\,6} = -\,\beta^4(\beta - 1728)^3\,
  \frac{\Delta(\Lambda_\idealc)}{\Delta(\Zi)},
\]
\PaperI[Prop.~6.10, eq.~(7.1)]: the phase is a weight-$2$ ratio at two CM
points, and its sixth power a $\Delta$-quotient dressed by singular moduli.
This paper takes that statement to its quantitative end.  The theme is
single: \emph{the phases of the Schmidt arrangement form an elliptic-unit
system, and the index of that system is a class number.}

\subsection*{The arc}
Write, for a class $\idealc$ of a level with reduced CM point
$\tau_\idealc$, $g_\idealc = y_\idealc^6|\Delta_q(\tau_\idealc)|$ for its
scale-invariant $\Delta$-datum, and, for a character $\chi$ of the class
group,
\[
  S(\chi) = \sum_\idealc\chi(\idealc)\log|u_\idealc|,
  \qquad
  \Sigma_x(\chi) = \sum_\idealc\chi(\idealc)\log|j(\idealc) - x| .
\]
Taking logarithms of
the closed forms and summing against $\chi$, the first Kronecker limit
formula converts $\sum\chi\log g$ into $-12L'(0,\chi)$ for the Epstein
$L$-function of the level (Proposition \ref{prop:klf}), and one obtains
the two \emph{master identities} (Theorems \ref{thm:euc-master} and
\ref{thm:hyp-master}):
\begin{gather*}
  \text{Euclidean:}\quad
  S(\chi) = -2\,L'(0,\chi) + \tfrac23\,\Sigma_0(\chi) + \tfrac12\,\Sigma_{1728}(\chi);\\
  \text{hyperbolic:}\quad
  S(\chi) = \begin{cases} 0, & \chi(\idealr_n) = +1,\\
  -4\,L'(0,\chi) + \tfrac43\,\Sigma_0(\chi) + \Sigma_{1728}(\chi), &
  \chi(\idealr_n) = -1.\end{cases}
\end{gather*}
The $L'$-part is the Kronecker limit formula; the $j$-parts are
Gross--Zagier collision quantities.  Two structural facts make the
hyperbolic identity possible: the Norm Lemma of \PaperI[Lemma~6.11] kills
both the unit $\eps = n + \sqrt{n^2-1}$ and the cocycle scalar $\mu$ in
$|u_f|$, leaving a scale-invariant two-CM-point law (Lemma
\ref{lem:hyp-classwise}); and the twist law $u_fu_{\idealr_nf} = 1$ of
\PaperI[Thm.~6.6] forces the character sums to live on the odd
characters.

For genus characters the identities evaluate completely (Section
\ref{sec:genus}).  The Epstein $L$-function of a real character factors
as a product of two Dirichlet $L$-functions times a finite Euler
correction at square conductors --- an identity of $q$-expansions, hence
provable by Sturm's bound --- and $L'(0,\chi)$ becomes an explicit
rational multiple of $\log\eps_{d_2}$ for a fundamental unit.  On the
Euclidean side the real quadratic field is $\Q(\sqrt n)$ in every case:
the phases of the curvature-$2n$ disks know the Pell unit of $\Q(\sqrt n)$.
The $j$-dressings $\Sigma_0, \Sigma_{1728}$ factor exactly in the same
field as unit powers times ratios of split-prime generators, supported on
the Gross--Zagier sets $\GZ(D,-3)$ and $\GZ(D,-4)$; for the non-real
characters certified integer-relation searches find no such expression,
as the Stark philosophy predicts.

The $\Delta$-part of the hyperbolic phase is a genuine unit.  With
$r_0 = \tfrac{n-1}2$ and $\idealr = \lat{r_0, \sqrt D/2}$ the invertible
ambiguous twist ideal of the order of discriminant $D = 1 - n^2$, the
$\idealr$-twisted $\Delta$-ratio
$R_f = u_f^{6}\,\beta_2^4(\beta_2-1728)^3/\beta_1^4(\beta_1-1728)^3$ has
the lattice form $R_f = r_0^{6}\Delta(\idealb_1)/\Delta(\idealr^{-1}\idealb_1)$
(Lemma \ref{lem:lattice}), and Theorem \ref{thm:unit} proves that $R_f$
is an algebraic unit for every class of every odd level, on every
imprimitive stratum as well.  The proof is local: the twist is an isogeny
whose tangent valuation is the same for every proper ideal (a rigidity of
$p$-divisible-group quotients), and composing two twists gives
multiplication by $p^k$, which balances the valuation exactly against
$r_0^6$.  No Kronecker limit formula, no Siegel units and no
quasi-canonical liftings enter.  The same rigidity gives the per-class
valuation law for the Euclidean $\Delta$-data (Theorem \ref{thm:perclass}):
$v_\idealP(n^{12}\Delta(\Lambda_\idealc)/\Delta(\Zi))$ is independent of
$\idealc$, with an explicit value $w_p(k)$, which settles the
split-prime denominator ladders of \PaperI[\S7] exactly and decomposes
the inert ones into certified single-slope ingredients.  A canonical
sixth root $w_f$ of $R_f$ obeys the mirror and twist laws at first power
and is an $\eta$-quotient on $\Gamma_0(r_0)$ at a Heegner point (Section
\ref{sec:firstpower}); at level $7$ the first-power Schmidt units are
$w = -(2 \pm \sqrt3)$, the fundamental unit of $\Q(\sqrt3)$ up to sign.

Finally, the index (Section \ref{sec:robert}).  On the quadratic layer,
the eigenprojections of the unit systems under a real character $\chi$
are explicit integer multiples of $\log\eps_{d_2}$, with coefficients
built from the class numbers $h(d_1)$ and $h(d_2)$.  On the
cubic layer --- the eight Euclidean levels $n \le 27$ whose ring class
group has $3$-torsion --- the principal coset product $\theta$ of the
$\Delta$-data, normalized by the cube root of the $\Delta$-mass, is a unit
$\theta_u$ of the real cubic subfield $L_3$ of the ring class field, and
\[
  \bigl[\Oc_{L_3}^\times : \langle -1, \theta_u\rangle\bigr]
  \;=\; 8\,h_{L_3}\,C_n(0)
\]
(Theorem \ref{thm:robert}), where $C_n(0)$ is the integer imprimitivity
multiplier, equal to $1$ at the primitive levels $9, 11, 13, 23$ where
$h_{L_3} = 1, 1, 3, 2$.  The identity $\text{index} = 8L'(0,\chi_3)/R_{L_3}$
is proved (Propositions \ref{prop:thetaunit}--\ref{prop:cubicfield}); the
content of the theorem is the certified computation of the fundamental
units, rigorous through a root descent capped by Friedman's regulator
bound.  The row $n = 23$ was computed after the law $8h_{L_3}$ had been
formulated on $9, 11, 13$, and confirmed it on a new cubic field of class
number $2$: \emph{the class number of the cubic field is the Robert index
of the Schmidt units}, the exact analogue of ``cyclotomic units have index
$h^+$'', realized by the disks of curvature $2n$.

\subsection*{Position in the literature}
The objects are classical and the ambition is to place a new geometric
invariant among them.  Kronecker's limit formula and its use to produce
units in class fields of imaginary quadratic fields go back to Kronecker
and Hecke, in the form we use to Siegel's lectures \cite{Siegel1961} and
Ramachandra \cite{Ramachandra1964}; Robert \cite{Robert1973} defined the
elliptic units of an abelian extension of an imaginary quadratic field and
proved that their index in the full unit group is the class number up to
an explicit factor, refined by Gillard--Robert \cite{GillardRobert1979}
and systematized, with the function-theoretic side of Siegel and Klein
forms, by Kubert--Lang \cite[Ch.~11--13]{KubertLang1981}; index formulas
for ramified elliptic units in the style of Sinnott are due to Oukhaba
\cite{Oukhaba2003}; Schertz's monograph \cite[Ch.~6--7]{Schertz2010}
presents the relation between elliptic units and class numbers and the
Stark units of ring class fields.  Stark \cite{Stark1980} showed that his
conjecture over an imaginary quadratic base is a theorem through
elliptic units; for ring class fields of prime conductor the Stark unit
generates, with the roots of unity, a subgroup of index exactly the class
number of the ring class field --- the form in which Hajir and
Rodriguez Villegas \cite{HajirVillegas1997} made the units explicit and in
which K\"u\c{c}\"uksakall{\i} \cite{Kucuksakalli2012} computed class numbers
of ring class fields.  On the finite-place side, Gross--Zagier's
factorization of singular moduli \cite{GrossZagier1985} is the reference
for the collision quantities $\Sigma_0, \Sigma_{1728}$, and Gross's
quasi-canonical liftings \cite{Gross1986} give the conceptual reading of
the valuation laws of Section \ref{sec:unit} --- a reading we keep as
interpretation, not input.  Meyer's monograph \cite{Meyer1957} computed
class numbers of abelian extensions of quadratic fields through Kronecker
limit formulas, and the higher Kronecker limit formulas of Vlasenko--Zagier
\cite{VlasenkoZagier2013} along real quadratic data are the closest
existing relatives of $\log|u_f|$ on the real-quadratic side, where the
phase is normalized by the unit of $\Q(\sqrt{n^2-1})$ and was shown in
\PaperI[Prop.~6.29] not to be cycle-integral data.  Friedman's regulator
bound \cite{Friedman1989} is the tool that makes the fundamental units of
Section \ref{sec:robert} rigorous.

What is new is not the Kronecker limit formula, nor the existence of
units in ring class fields with class-number index --- those are the
classical theory just cited.  It is that a \emph{geometric} invariant of
circles is such a system: the character sums of the phase compute
$L'(0,\chi)$ exactly, with the unit $\eps$ cancelling and the genus field
$\Q(\sqrt n)$ appearing; the ambiguous-ideal twist $\idealr$ of a
non-maximal order preserves every local $\Delta$-datum, so that the
$\Delta$-ratios the phase carries are units on every stratum; the
Euclidean $\Delta$-data have class-independent valuations; and the
$\Delta$-coset units of the curvature-$2n$ disks have index
$8h_{L_3}C_n(0)$ in the cubic layer over $\Q(i)$, a number computed and
certified at eight levels.  A literature record accompanies the paper in
the repository; we found no prior index computation for elliptic-unit
systems in ring class fields of $\Q(i)$ and no prior appearance of these
$\Delta$-coset units.  How $\theta_u$ compares with the Stark unit of the
ring class field --- both generate finite-index subgroups, with indices
$8h_{L_3}$ in $\Oc_{L_3}^\times$ and $h_{H_n}$ in $\Oc_{H_n}^\times$
respectively --- is a natural question we do not pursue here.

\subsection*{What is proved and what is certified}
Proved for every level: the limit-value formula (Proposition
\ref{prop:klf}); the classwise laws (Lemmas \ref{lem:euc-classwise},
\ref{lem:hyp-classwise}); both master identities (Theorems
\ref{thm:euc-master}, \ref{thm:hyp-master}) with the trivial-character
anchor; the $\Delta$-mass polynomial and the uniform Stark law (Theorem
\ref{thm:Dn}); the rationality of the twisted-ratio polynomial (Theorem
\ref{thm:Rpoly}); the genus factorizations and the resulting $L'(0,\chi_2)$
closed forms at every displayed level (Proposition \ref{prop:genus-euc},
Sturm's bound); each displayed $\Sigma$-factorization (exact arithmetic in
its quadratic field); the lattice lemma and the unit theorem on all strata
(Lemma \ref{lem:lattice}, Theorem \ref{thm:unit}); the collapse $R \equiv 1$
on strata with principal induced twist; the per-class valuation law and
the split ladder (Theorem \ref{thm:perclass}, Corollary \ref{cor:split});
the first-power laws of $w_f$ and its $\eta$-quotient form (Proposition
\ref{prop:wlaws}); the unit property of $\theta_u$, the Stark relation and
the index identity $8L'(0,\chi_3)/R_{L_3}$ (Propositions
\ref{prop:thetaunit}--\ref{prop:cubicfield}).  Certified, by exact
integer arithmetic or with at least $140$ (typically $\ge 200$) spare
digits, and not proved in general: the coefficient values of the displayed
polynomials and their irreducibility; the Newton-polygon verification of
Theorem \ref{thm:perclass} and the inert-ladder decompositions; the
coherence table $m(n)$ of the first-power units; the fundamental units,
class numbers and indices of Theorem \ref{thm:robert}; the
$\eps$-ratio identifications at the composite conductor $n = 15$; and the
converse direction of the stratum criterion.  Certified negatives: the
PSLQ non-fits for non-real characters and the absence of a congruence
law for $m(n)$.

\subsection*{Plan of the paper}
Section \ref{sec:master} proves the limit-value formula, the classwise
laws and the master identities, and records the $\Delta$-mass and
twisted-ratio polynomials.  Section \ref{sec:genus} evaluates the genus
characters in closed form.  Section \ref{sec:unit} proves the unit theorem
and the per-class valuation law --- the technical core.  Section
\ref{sec:firstpower} constructs the first-power Schmidt units.  Section
\ref{sec:robert} computes the Robert index.  Section \ref{sec:outlook}
lists the four genuine continuations, and Appendix \ref{app:verification}
documents the machine verification.

\subsection*{Status conventions}
As in \PaperI, a \textbf{Theorem}, \textbf{Proposition}, \textbf{Lemma} or
\textbf{Corollary} without qualification is proved.  A statement labelled
\emph{(computational)} is a finite claim verified by exact arithmetic, or
a numerical claim certified at a stated precision with stated safety
margins, by a named script; the range of verification is part of the
statement.  Every script cited was re-run for this paper (Appendix
\ref{app:verification}).  No integer-relation search is used to establish
a claim: the discoveries made by PSLQ are re-derived in exact arithmetic,
and PSLQ appears only in certified \emph{negative} statements.

\subsection*{Notation}
$K$ denotes the imaginary quadratic field of the aspect at hand.
\emph{Euclidean level $n \ge 2$}: $K = \Q(i)$, $\Oc_n = \Z + n\Zi$ the
order of conductor $n$, discriminant $D = -4n^2$, class group
$\Cl(\Oc_n)$ of order $h = h(-4n^2)$, ring class field $H_n$; the phases
$u_\idealc = \Theta/\Omega$ of \PaperI[\S7], normalized by the lemniscatic
period, with $|u_\idealc|$ and $u_\idealc^2$ well defined on classes
\PaperI[Prop.~7.10]; $\Lambda_\idealc \subseteq \Zi$ the primitive
index-$n$ sublattice of the class \PaperI[Thm.~7.6];
$G_\idealc := n^{12}\Delta(\Lambda_\idealc)/\Delta(\Zi)$ and the
$\Delta$-mass $M(n) = \prod_\idealc G_\idealc$ \PaperI[Thm.~7.18].
\emph{Hyperbolic level $n \ge 3$ odd}: $K = \Q(\sqrt{1-n^2})$,
$D = 1 - n^2$, $N = n^2 - 1$, $\Oc$ the order of discriminant $D$ with
class group $\Cl(D)$ of primitive classes, $h = h(D)$, ring class field
$H$; $\eps = n + \sqrt N$; $r_0 = \tfrac{n-1}2$, $s_0 = \tfrac{n+1}2$,
$\omega_0 = \sqrt D/2$; the twist ideal $\idealr = \idealr_n =
\lat{r_0, \omega_0}$, invertible and ambiguous, with $\bar\idealr =
\idealr$, $\idealr^2 = (r_0)$, $N\idealr = r_0$, and class
$[\idealr_n] = [(r_0, 0, s_0)]$ \PaperI[Thm.~4.5]; for a primitive class
$f$ with CM point $m_1$, $\idealb_1 = \lat{1, m_1}$, $\beta_1 = j(f)$,
$\beta_2 = j(\idealr f) = j(\idealr^{-1}\idealb_1)$ (the complex conjugate
of the moduli coordinate $\beta_2$ of \PaperI[\S6.1]), and $u_f = \eps\Theta_f$
the normalized phase with laws $u_{f^{-1}} = \bar u_f$ and
$u_fu_{\idealr f} = 1$ \PaperI[Thm.~6.6].  In both aspects the unit count of
the order is $w = 2$.  $\Delta_q = \eta^{24}$ is the normalized
discriminant, $\Delta(\Lambda)$ the weight-$12$ lattice discriminant,
$j$ the modular invariant, $\tau_\idealc$ the CM point of the reduced
form $Q_\idealc = (A, B, C)$ of a class, $y_\idealc = \im\tau_\idealc =
\sqrt{|D|}/2A$, and
\begin{equation}\label{eq:gdef}
  g_\idealc \;:=\; y_\idealc^{\,6}\,\bigl|\Delta_q(\tau_\idealc)\bigr|
\end{equation}
the scale-invariant $\Delta$-datum of the class ($g = \covol^6|\Delta|$ of
any lattice in the class, up to one global constant).  For a real
quadratic discriminant $d$, $\eps_d > 1$ is the fundamental unit and
$h(d)$ the class number; $w(d)$ is the number of roots of unity of the
imaginary quadratic order of discriminant $d < 0$.

\subsection*{Acknowledgements}
The results of this paper were obtained in a sequence of
computer-assisted research sessions; the complete research notes,
verification scripts and literature record are available in the
accompanying repository.
%% TODO before submission: repository URL / archival DOI.

%%% ============================================================
\section{The classwise laws and the master identities}\label{sec:master}

\subsection{The Epstein $L$-function of a level and the limit value}
\label{subsec:epstein}

\begin{definition}\label{def:epstein}
For a discriminant $D < 0$ with class group $\Cl(D)$ of primitive forms
and a character $\chi$ of $\Cl(D)$, set
\[
  L(s,\chi) \;:=\; \frac1w\sum_{\idealc \in \Cl(D)}\chi(\idealc)\,
  \zeta_{Q_\idealc}(s),
  \qquad
  \zeta_Q(s) = \sideset{}{'}\sum_{(x,y)\in\Z^2}Q(x,y)^{-s},
\]
the form-class (Epstein) $L$-function of the discriminant.  For $m \ge
1$ let $r_\idealc(m)$ be the number of representations of $m$ by
$Q_\idealc$ and $R_\chi(m) := \sum_\idealc\chi(\idealc)r_\idealc(m)$, so
that $L(s,\chi) = w^{-1}\sum_{m \ge 1}R_\chi(m)m^{-s}$.
\end{definition}

Since $(x,y) \mapsto (x,-y)$ bijects the representations by $Q$ and by
$Q^{-1}$, one has $R_\chi = R_{\bar\chi}$, and $L(s,\chi) = L(s,\bar\chi)$
is real on the real axis.  In our two aspects $D = -4n^2$ resp.\ $D = 1
- n^2$, and $w = 2$.

\begin{proposition}[functional equation and the limit value]\label{prop:klf}
Write $\gamma(s) = \bigl(\tfrac{\sqrt{|D|}}{2\pi}\bigr)^{s}\Gamma(s)$.
Then $\gamma(s)\zeta_Q(s) = \gamma(1-s)\zeta_Q(1-s)$ for every class, so
for nontrivial $\chi$ the function $L(s,\chi)$ is entire with $L(0,\chi)
= 0$, and
\begin{equation}\label{eq:klf}
  L'(0,\chi) \;=\; \frac{\sqrt{|D|}}{2\pi}\,L(1,\chi)
  \;=\; -\,\frac{1}{12}\sum_{\idealc}\chi(\idealc)\,\log g_\idealc .
\end{equation}
\end{proposition}

\begin{proof}
With $\tau_Q = \tfrac{-B + i\sqrt{|D|}}{2A}$ one has $A|x + y\tau_Q|^2 =
Q(x,-y)$ and $Ay_{\tau_Q} = \sqrt{|D|}/2$, whence
$\zeta_Q(s) = \bigl(\tfrac{2}{\sqrt{|D|}}\bigr)^{s}E(\tau_Q, s)$ for the
nonholomorphic Eisenstein series $E(\tau,s) = \sum'y^s/|m\tau + n|^{2s}$.
The completed series $E^*(\tau,s) = \pi^{-s}\Gamma(s)E(\tau,s)$ satisfies
$E^*(\tau,s) = E^*(\tau,1-s)$, and $\gamma(s)\zeta_Q(s) = E^*(\tau_Q,s)$:
this is the functional equation, class by class.  For nontrivial $\chi$
the poles at $s = 1$ cancel in the character sum, so $\Gamma(s)L(s,\chi)$
is finite at $s = 0$; since $\Gamma(s) \sim 1/s$, $L(0,\chi) = 0$ and
$L'(0,\chi) = \lim_{s \to 0}\gamma(s)L(s,\chi) = \gamma(1)L(1,\chi)$.  The
first Kronecker limit formula \cite[Ch.~I]{Siegel1961}, \cite[Ch.~20]{Lang1987},
\[
  E(\tau,s) = \frac{\pi}{s-1} + 2\pi\bigl(\gamma_E - \log 2 -
  \log(\sqrt y\,|\eta(\tau)|^2)\bigr) + O(s-1),
\]
gives, after the $\chi$-sum has killed every $\idealc$-independent
constant, and with $\sqrt y\,|\eta|^2 = (y^6|\Delta_q|)^{1/12} =
g^{1/12}$,
\[
  L(1,\chi) = -\frac1w\cdot\frac{2}{\sqrt{|D|}}\cdot\frac{\pi}{6}
  \sum_\idealc\chi(\idealc)\log g_\idealc ;
\]
multiply by $\gamma(1) = \sqrt{|D|}/2\pi$ and use $w = 2$.
\end{proof}

The verification of every identity below rests on an evaluation of
$L'(0,\chi)$ that uses \emph{no} modular quantity at all.

\begin{lemma}[independent evaluation]\label{lem:gamma-eval}
For nontrivial $\chi$, with $\alpha = 2\pi/\sqrt{|D|}$ and $E_1$ the
exponential integral,
\[
  L'(0,\chi) \;=\; \frac1w\sum_{m \ge 1}R_\chi(m)\Bigl[\frac{e^{-\alpha m}}
  {\alpha m} + E_1(\alpha m)\Bigr] .
\]
\end{lemma}

\begin{proof}
Let $\theta_Q(t) = \sum_{(x,y)\in\Z^2}e^{-\alpha Q(x,y)t}$.  The Mellin
transform gives
\[
  \gamma(s)\zeta_Q(s) = \int_0^\infty(\theta_Q(t) - 1)\,t^{s-1}\,dt,
\]
and Poisson summation gives $\theta_Q(1/t) = t\,
\theta_{Q^{-1}}(t)$: in the normalization $\alpha = 2\pi/\sqrt{|D|}$ the
Gram matrix $2G_Q/\sqrt{|D|}$ has determinant $1$, and its inverse is
$2G_{Q^{-1}}/\sqrt{|D|}$.  Splitting the integral at $t = 1$ and
inverting $t$ on $(0,1)$,
\[
  \gamma(s)\zeta_Q(s) = \int_1^\infty(\theta_Q(t) - 1)t^{s-1}\,dt
  + \int_1^\infty(\theta_{Q^{-1}}(t) - 1)t^{-s}\,dt + \frac1{s-1} -
  \frac1s .
\]
Summing against a nontrivial $\chi$ kills the last two terms, and
$R_\chi = R_{\bar\chi}$ identifies the $Q^{-1}$-sum with the $Q$-sum;
at $s = 0$ the left side is $L'(0,\chi)$ by Proposition \ref{prop:klf},
and the two integrals of $e^{-\alpha mt}$ are $E_1(\alpha m)$ and
$e^{-\alpha m}/\alpha m$.
\end{proof}

The series converges like $e^{-\alpha m}$; at $250$ digits it needs
between $600$ and $2600$ terms at the levels below, and its inputs are
the exact integers $R_\chi(m)$.  Every limit-formula identity in this
paper is checked against it (residuals $\le 2\cdot 10^{-248}$ at $250$
digits; Appendix \ref{app:verification}).

\subsection{Euclidean: the classwise law and the master identity}
\label{subsec:euc-master}

\begin{lemma}[classwise law]\label{lem:euc-classwise}
For every class $\idealc \in \Cl(\Oc_n)$ of every level $n \ge 2$, with
$\beta = j(\idealc)$,
\[
  |u_\idealc|^{6} \;=\; \bigl|\beta^4(\beta - 1728)^3\bigr|\;
  \frac{g_\idealc}{n^{6}\,|\Delta_q(i)|} .
\]
\end{lemma}

\begin{proof}
$u^6 = -\beta^4(\beta - 1728)^3\Delta(\Lambda_\idealc)/\Delta(\Zi)$ is
\PaperI[eq.~(7.1)].  Write $\Lambda_\idealc = c\lat{1, w}$; then $|c|^2y_w
= n = \covol(\Lambda_\idealc)$ and $|\Delta(\Lambda_\idealc)| =
|c|^{-12}(2\pi)^{12}|\Delta_q(w)|$, so $|\Delta(\Lambda_\idealc)/
\Delta(\Zi)| = y_w^6|\Delta_q(w)|/(n^6|\Delta_q(i)|)$, and $y_w^6
|\Delta_q(w)| = g_\idealc$ by $\SL_2(\Z)$-invariance of $y^6|\Delta_q|$.
\end{proof}

\begin{theorem}[Euclidean master identity]\label{thm:euc-master}
For every nontrivial character $\chi$ of $\Cl(\Oc_n)$, with $S(\chi) =
\sum_\idealc\chi(\idealc)\log|u_\idealc|$ and $\Sigma_x(\chi) =
\sum_\idealc\chi(\idealc)\log|j(\idealc) - x|$,
\begin{equation}\label{eq:euc-master}
  \boxed{\;S(\chi) \;=\; -2\,L'(0,\chi) + \tfrac23\,\Sigma_0(\chi) +
  \tfrac12\,\Sigma_{1728}(\chi).\;}
\end{equation}
For the trivial character, with $H_{-4n^2}$ the ring class polynomial
and $M(n)$ the $\Delta$-mass of \PaperI[Thm.~7.18],
\begin{equation}\label{eq:euc-trivial}
  6\sum_\idealc\log|u_\idealc| \;=\; 4\log|H_{-4n^2}(0)| +
  3\log|H_{-4n^2}(1728)| + \log|M(n)| - 12\,h\log n .
\end{equation}
\end{theorem}

\begin{proof}
Take logarithms in Lemma \ref{lem:euc-classwise} and sum against
$\chi$: the $\idealc$-independent constant $-6\log n - \log|\Delta_q(i)|$
dies for $\chi \ne 1$, and $\sum\chi\log g = -12L'(0,\chi)$ by
Proposition \ref{prop:klf}.  For the trivial character sum the same
identity over all classes: $\sum\log|\beta| = \log|H(0)|$, $\sum\log|\beta
- 1728| = \log|H(1728)|$, and $\sum_\idealc(\log g_\idealc -
\log|\Delta_q(i)| - 6\log n) = \log|M(n)| - 12h\log n$ by the definition
$M(n) = \prod_\idealc n^{12}\Delta(\Lambda_\idealc)/\Delta(\Zi)$.
\end{proof}

Both statements are verified at $n = 3, 5, 7, 9, 11, 13, 15$, every
character, with residuals $\le 2\cdot 10^{-248}$ ($250$ digits; $400$ at
$n = 15$), the identity and the independent evaluation of Lemma
\ref{lem:gamma-eval} being separate computations; the trivial-character
anchor is checked against the certified $H_{-4n^2}$ and the closed form of
$M(n)$.

The limit formula becomes an algebraic statement through one integer
polynomial per level.

\begin{theorem}[the $\Delta$-mass polynomial and the uniform Stark law]
\label{thm:Dn}
Let $G_\idealc = n^{12}\Delta(\Lambda_\idealc)/\Delta(\Zi)$; this is a
nonzero algebraic integer, of absolute value $n^6g_\idealc/|\Delta_q(i)|$.
Then:
\begin{enumerate}
\item $\sigma(G_\idealc) = G_{\idealc(\sigma)^{-1}\idealc}$ for $\sigma \in
\Gal(\bar\Q/K)$ with Artin class $\idealc(\sigma)$ on $H_n$, and
$\overline{G_\idealc} = G_{\idealc^{-1}}$; hence the multiset
$\{G_\idealc\}$ is $\Gal(\bar\Q/\Q)$-stable, and
\[
  D_n(x) := \prod_{\idealc \in \Cl(\Oc_n)}(x - G_\idealc) \;\in\; \Z[x],
  \qquad D_n(0) = (-1)^h M(n) .
\]
\item For every nontrivial $\chi$:
$\;-12\,L'(0,\chi) = \sum_\idealc\chi(\idealc)\log|G_\idealc|$.
\end{enumerate}
Thus all the $L'(0,\chi)$ of a level are character combinations of the
logarithms of the roots of one integer polynomial: the $\Delta$-data
$\{G_\idealc\}$ of the curvature-$2n$ disks form an elliptic-unit system
for the ring class tower of $\Q(i)$ in the sense of Siegel and Robert.
\end{theorem}

\begin{proof}
(1) $G_\idealc = F(\Lambda_\idealc; \Zi)$ with $F(\Lambda; L) =
n^{12}\Delta(\Lambda)/\Delta(L)$, isobaric of joint weight zero.  The
transport lemma \PaperI[Lemma~7.13] gives $\sigma(F(\Lambda_\idealc;\Zi))
= F(s^{-1}\Lambda_\idealc; s^{-1}\Zi)$ for the id\`ele $s$ of $\sigma$;
rescaling jointly by a generator of the principal ideal
$(s^{-1}\Zi)^{-1}$, exactly as in the translation law \PaperI[Prop.~7.14],
lands on a primitive representative of the class
$\idealc(\sigma)^{-1}\idealc$, determined up to $\mu_4$ --- and
$\Delta(i\Lambda) = \Delta(\Lambda)$, so there is no ambiguity at all.
Complex conjugation: $\Delta(\bar\Lambda) = \overline{\Delta(\Lambda)}$
and $\bar\Lambda$ represents $\idealc^{-1}$ \PaperI[Prop.~7.10].  The
coefficients of $D_n$ are therefore rational, and integral because
$n^{12}\Delta$-quotients of commensurable CM lattices are algebraic
integers \cite[Ch.~12]{Lang1987}; $D_n(0) = (-1)^h\prod G_\idealc =
(-1)^hM(n)$.  (2) $\log|G_\idealc| = \log g_\idealc + 6\log n -
\log|\Delta_q(i)|$, so the character sum reproduces Proposition
\ref{prop:klf}.
\end{proof}

The certified polynomials (integer coefficients with $\ge 226$ spare
digits at $250$ digits, $358$ at $n = 15$; constant terms equal to
$(-1)^hM(n)$ exactly):
\begin{align*}
D_3 &= x^2 - 378x + 729, \qquad D_5 = x^2 - 322x + 1,\\
D_7 &= x^4 + 58604x^3 + 5618502582x^2 - 191605917748x + 7^{6},\\
D_9 &= x^6 + 777114x^5 + 469570343535x^4 + 88033943990172204x^3\\
    &\qquad + 4975970841324612804303x^2 - 6464018128844293554150x + 3^{24},\\
D_{11} &= x^6 - 2940366x^5 + 2977596073815x^4 + 116499046538184860x^3\\
    &\qquad + 38813935158542440699935x^2 - 1046289507684469650030x + 11^{6},\\
D_{13} &= x^6 + 7210938x^5 + 12700046391567x^4 - 1529217345755339156x^3\\
    &\qquad + 61048319821249936271631x^2 - 22813565655771828294x + 1,\\
D_{15} &= x^8 - 243432x^7 - 49497426055620x^6 - 1528089585117383803032x^5\\
    &\qquad + 12625658759774046395283174x^4\\
    &\qquad + 2831276695986378100679863409676072x^3\\
    &\qquad + 85367425855325247262124045173735929698940x^2\\
    &\qquad - 331767481864968679465735138096222248x + 3^{24}.
\end{align*}

\begin{remark}[coset cubics; the Stark objects of the cubic layer]
\label{rem:cosetcubics}
At $n = 9, 11, 13$, where $\Cl(\Oc_n) \cong \Z/6$, the products of the
$G_\idealc$ over the three cosets of the kernel of the order-$3$
character are the roots of the certified integer cubics
\begin{align*}
n = 9:&\quad x^3 - 76822296291x^2 + 4976085199686658256835x - 3^{24},\\
n = 11:&\quad x^3 + 277946532501x^2 + 38813938298955894507411x - 11^{6},\\
n = 13:&\quad x^3 + 446643445245x^2 + 61048319249786206560771x - 1,
\end{align*}
with the Stark-type closed form
$-12L'(0,\chi_3) = \tfrac32\log|G^{(3)}_0| - \tfrac12\log|C_\Delta(0)|$,
$G^{(3)}_0$ the principal coset product and $C_\Delta(0)$ the constant
term, verified to $10^{-247}$.  These cubics are the objects whose unit
normalization carries the Robert index of Section \ref{sec:robert}.
\end{remark}

\subsection{Hyperbolic: cancellation of $\eps$ and $\mu$, and the
odd-character law}\label{subsec:hyp-master}

\begin{lemma}[scale-invariant closed form]\label{lem:hyp-classwise}
For every primitive class $f$ of discriminant $D = 1 - n^2$, with
$\beta_1 = j(f)$, $\beta_2 = j(\idealr_nf)$ and $g(f) =
y_f^6|\Delta_q(\tau_f)|$,
\[
  |u_f|^{6} \;=\;
  \frac{\bigl|\beta_1^4(\beta_1 - 1728)^3\bigr|}
       {\bigl|\beta_2^4(\beta_2 - 1728)^3\bigr|}\cdot
  \frac{g(f)}{g(\idealr_nf)} .
\]
In particular neither $\eps = n + \sqrt{n^2-1}$ nor $\mu$ survives in
$|u_f|$.
\end{lemma}

\begin{proof}
By the closed form \PaperI[Prop.~6.10], $u_f = -\eps\mu^{-2}h_2(\idealb_1)
/h_2(\idealb_2)$ with $\idealb_2 = \lat{1, -m_2}$ the lattice of the class
$[\idealr_nf]$ \PaperI[Prop.~6.7], and by the Norm Lemma
\PaperI[Lemma~6.11], $|\mu|^2 = \eps q_1/q_2$ with $q_1, q_2$ the
half-curvatures of the two circles; hence $|u_f| =
(q_2/q_1)\,|h_2(\idealb_1)/h_2(\idealb_2)|$.  The weight algebra $h_2^6 =
c\,j^4(j - 1728)^3\Delta$ (one universal constant $c$) and
\[
  |\Delta(\idealb_i)| = (2\pi)^{12}\,\frac{g(\idealb_i)}{\covol(\idealb_i)^6},
  \qquad
  \frac{\covol(\idealb_1)}{\covol(\idealb_2)} = \frac{q_2}{q_1}
\]
(both CM points have imaginary part $\sqrt N/2q_i$) give
\[
  \Bigl|\frac{h_2(\idealb_1)}{h_2(\idealb_2)}\Bigr|^{6}
  = \Bigl(\frac{q_1}{q_2}\Bigr)^{6}\cdot
  \frac{|\beta_1^4(\beta_1-1728)^3|}{|\beta_2^4(\beta_2-1728)^3|}\cdot
  \frac{g_1}{g_2} ;
\]
assemble, and identify $\beta_2, g_2$ as the class functions at
$\idealr_nf$ through $[\idealb_2] = [\idealr_n][\idealb_1]$
\PaperI[Thm.~4.5].
\end{proof}

The lemma is verified classwise at $n = 9, 11, 13, 15$ with maximal
deviation $4.5\cdot 10^{-249}$.  The expected new ingredient of the
hyperbolic aspect, $\log\eps_{n^2-1}$ coming from $\eps^6\mu^{-12}$,
cancels identically: the Norm Lemma places $\eps$ exactly as the
discrepancy between the two archimedean places of $\mu$, and $|u_f|$ sees
only one of them.

\begin{theorem}[hyperbolic master identity]\label{thm:hyp-master}
Let $\chi$ be a character of $\Cl(1 - n^2)$, $S(\chi) = \sum_f\chi(f)
\log|u_f|$, $\Sigma_x(\chi) = \sum_f\chi(f)\log|\beta_1(f) - x|$.  Then
\begin{equation}\label{eq:hyp-master}
  \boxed{\;S(\chi) \;=\; \begin{cases} 0, & \chi(\idealr_n) = +1,\\[3pt]
  -4\,L'(0,\chi) + \tfrac43\,\Sigma_0(\chi) + \Sigma_{1728}(\chi), &
  \chi(\idealr_n) = -1 .\end{cases}\;}
\end{equation}
\end{theorem}

\begin{proof}
Sum $6\log|u_f|$ from Lemma \ref{lem:hyp-classwise} against $\chi$.
Re-indexing the $\idealr_nf$-terms by $f \mapsto \idealr_nf$ multiplies
their sum by $\chi(\idealr_n)^{-1} = \chi(\idealr_n)$, the twist class
being $2$-torsion, so
$6S(\chi) = (1 - \chi(\idealr_n))\bigl[4\Sigma_0 + 3\Sigma_{1728} +
\Sigma_\gamma\bigr]$ with $\Sigma_\gamma(\chi) = \sum_f\chi(f)\log g(f) =
-12L'(0,\chi)$ by Proposition \ref{prop:klf} at discriminant $1 - n^2$.
\end{proof}

The vanishing statement is the twist law $u_{\idealr f}u_f = 1$ seen
through the limit formula; both branches are verified at $n = 9, 11, 13,
15$, every character, with residuals $\le 10^{-248}$, the even-character
sums vanishing to $10^{-249}$ or exactly.

\begin{theorem}[the twisted $\Delta$-ratio polynomial]\label{thm:Rpoly}
Set
\begin{equation}\label{eq:Rdef}
  R_f \;:=\; u_f^{6}\;\frac{\beta_2^4(\beta_2 - 1728)^3}
  {\beta_1^4(\beta_1 - 1728)^3},
  \qquad |R_f| = \frac{g(f)}{g(\idealr_nf)},
\end{equation}
a class function, with $R_{\idealr_nf} = R_f^{-1}$.  Then the multiset
$\{R_f\}_{f \in \Cl(D)}$ is stable under $\Gal(\bar\Q/\Q)$, so that
$\prod_f(x - R_f) \in \Q[x]$; and for every odd character
($\chi(\idealr_n) = -1$),
\begin{equation}\label{eq:R-stark}
  -24\,L'(0,\chi) \;=\; \sum_f\chi(f)\log|R_f| .
\end{equation}
\end{theorem}

\begin{proof}
$|R_f| = g(f)/g(\idealr_nf)$ is Lemma \ref{lem:hyp-classwise}, and
$R_{\idealr f} = 1/R_f$ is the twist law together with $\idealr^2 = 1$ in
$\Cl(D)$.  By the first-power dihedral law \PaperI[Thm.~6.19] and the
reciprocity for singular moduli, $u_f$, $\beta_1$ and $\beta_2$ are
transported by the same generalized dihedral permutation $f \mapsto
f^{e(\sigma)}\idealc(\sigma)$ of the classes (the translated pair is
again the canonical pair of the translated class), so $\sigma(R_f) =
R_{f^{e(\sigma)}\idealc(\sigma)}$.  Finally
$\sum_f\chi(f)[\log g(f) - \log g(\idealr_nf)] = (1 - \chi(\idealr_n))
\Sigma_\gamma(\chi) = 2\Sigma_\gamma(\chi) = -24L'(0,\chi)$ for odd $\chi$.
\end{proof}

This is the hyperbolic counterpart of Theorem \ref{thm:Dn}, and sharper:
the certified polynomials $\prod_f(x - R_f)$ (Section
\ref{subsec:unitrecord}) have \emph{integer} coefficients, are
palindromic, and have constant term $1$ --- the $R_f$ are algebraic units,
inverted in pairs by the $\idealr_n$-twist, so the $L'(0,\chi)$ of the odd
characters are $\chi$-combinations of logarithms of genuine units, with
no $S$-integer dressing at all in the $\Delta$-part.  Theorem
\ref{thm:unit} proves this for every $n$.

\subsection{The certified record}\label{subsec:record}

Character sums and $L'$-values at the computed levels ($20$ digits shown;
computed at $250$ digits, $400$ at Euclidean $n = 15$; $S(\chi) =
S(\bar\chi)$, one value per conjugate pair; characters of a non-cyclic
class group are indexed by their exponents on the generators of
$\Z/4\times\Z/2$).  Euclidean:
\begin{center}
\small
\begin{tabular}{@{}rlll@{}}
\toprule
$n$ & $\operatorname{ord}\chi$ & $S(\chi)$ & $L'(0,\chi)$\\
\midrule
3 & 2 & 10.119185642834999673 & 0.43898596564160556954\\
5 & 2 & 16.401110452980199012 & 0.96242365011920689500\\
7 & 2 & 50.898733868194595161 & 2.7686593833135738327\\
7 & 4 & 22.684295755688513572 & 1.4860221248769270925\\
9 & 2 & 18.137584684821439049 & 1.7559438625664222782\\
9 & 3 & 70.144229691414912554 & 4.0476417264498781634\\
9 & 6 & 34.382429251891407159 & 2.1364594185869239177\\
11 & 2 & 50.213109260828194576 & 2.9932228461263808979\\
11 & 3 & 86.033492426467556694 & 5.3026856597888262543\\
11 & 6 & 27.769444924657409111 & 2.3032180902028192162\\
13 & 2 & 23.802362926332249068 & 2.3895264345742186082\\
13 & 3 & 101.24102291036467905 & 6.5582440777920829100\\
13 & 6 & 50.399596052674761580 & 3.3904644593763457318\\
15 & 2 $(0,1)$ & 62.742974077314449734 & 3.8496946004768275800\\
15 & 2 $(2,0)$ & 123.17609219442540179 & 8.2537482755822421869\\
15 & 2 $(2,1)$ & 48.769335467586396360 & 3.5118877251328445563\\
15 & 4 $(1,0)$ & 132.06845592069919001 & 8.2852528669601990706\\
15 & 4 $(1,1)$ & 39.877919196363264725 & 3.4800432906097017186\\
\bottomrule
\end{tabular}
\end{center}
Hyperbolic (odd characters; the even ones have $S = 0$ to $10^{-249}$
or exactly):
\begin{center}
\small
\begin{tabular}{@{}rlll@{}}
\toprule
$n$ & $\chi$ & $S(\chi)$ & $L'(0,\chi)$\\
\midrule
9 & order 4 & 45.337283759443440194 & 1.0612750619050356520\\
11 & $(0,1)$ & 71.156367508225883636 & 1.7627471740390860505\\
11 & $(1,0)$ & 46.591472482040984242 & 1.2122976394880445490\\
13 & $(0,1)$ & 86.596498648832811158 & 2.2924316695611776878\\
13 & $(1,0)$ & 57.828173780263506754 & 1.5667992369724110787\\
15 & $(0,1)$ & 96.655012640029353845 & 2.7686593833135738327\\
15 & $(2,1)$ & 42.411391969797829839 & 1.7000422070566697504\\
15 & order 4 & 101.30861654091162593 & 2.5318972768037898006\\
\bottomrule
\end{tabular}
\end{center}
Two coincidences are visible to the eye and explained in Section
\ref{sec:genus}: Euclidean $n = 7$ (order $2$) and hyperbolic $n = 15$,
$\chi_{(0,1)}$, share $L' = \log(8 + 3\sqrt7)$, the Kronecker-symbol
pairs $(-7)\cdot 28$ and $(-8)\cdot 28$ arriving at the same
$\eps_{28}$; and the three real characters at Euclidean $n = 15$ have
$L' = 4\log\eps_{60}$, $8\log\eps_5$ and $\tfrac83\log\eps_{12}$.

%%% ============================================================
\section{Genus characters in closed form}\label{sec:genus}

For a real (genus) character every term of the master identities
evaluates in closed form: the $L'$-part through a factorization of the
Epstein $L$-function into two Dirichlet $L$-functions, the $j$-parts
through exact factorizations in a real quadratic field.

\subsection{Euclidean: the $L'$-part and the field $\Q(\sqrt n)$}
\label{subsec:genus-euc}

$\Cl(\Oc_n)$ is cyclic of order $6$ at $n = 9, 11, 13$, of order $2$ at
$n = 3, 5$, $\Z/4$ at $n = 7$, and $\Z/4\times\Z/2$ at $n = 15$; so there
is one nontrivial real character $\chi_2$ per level $n \le 13$ and three
at $n = 15$.  For a discriminant $d$ write $\chi_d = (\tfrac{d}{\cdot})$
for the Kronecker symbol, a (possibly imprimitive) Dirichlet character,
$L(s,\chi_d)$ for its $L$-function, and $d^*$ for the fundamental
discriminant of the primitive character inducing $\chi_d$.

\begin{proposition}[genus factorization of the Epstein $L$-function]
\label{prop:genus-euc}
For the real character $\chi_2$ of $\Cl(\Oc_n)$, $n \in \{3, 5, 7, 9, 11,
13\}$, and for the character $\chi_{(2,0)}$ at $n = 15$, the
representation numbers satisfy, for every $m \ge 1$,
\begin{equation}\label{eq:conv}
  R_{\chi_2}(m) \;=\; 2\,\bigl(\mathrm{conv}_{d_1,d_2} * \mathrm{corr}\bigr)(m),
  \qquad
  \mathrm{conv}_{d_1,d_2}(m) = \sum_{e \mid m}\chi_{d_1}(e)\,\chi_{d_2}(m/e),
\end{equation}
with the decompositions $d_1d_2 = -4n^2$ and finite corrections
$\mathrm{corr}$ (Dirichlet convolution; $\mathrm{corr} = \delta_1$ unless
stated) of the table below.  Equivalently, as Dirichlet series,
\[
  L(s,\chi_2) = L(s,\chi_{d_1})\,L(s,\chi_{d_2})\,C(s),
  \qquad C(s) = \sum_m\mathrm{corr}(m)m^{-s},
\]
and consequently
\begin{equation}\label{eq:Lprime-genus}
  L'(0,\chi_2) \;=\; \frac{2h(d_1^*)}{w(d_1^*)}\;h(d_2)\,\log\eps_{d_2}\;C(0).
\end{equation}
\begin{center}
\small
\begin{tabular}{@{}rlll@{}}
\toprule
$n$ & $d_1\cdot d_2 = -4n^2$ & correction $C(s)$ & $L'(0,\chi_2)$\\
\midrule
3 & $(-3)\cdot 12$ & --- & $\tfrac13\log\eps_{12}$\\
5 & $(-20)\cdot 5$ & --- & $2\log\eps_5$\\
7 & $(-7)\cdot 28$ & --- & $\log\eps_{28}$\\
9 & $(-27)\cdot 12$ & $1 + 3^{1-2s}$ & $\tfrac43\log\eps_{12}$\\
11 & $(-11)\cdot 44$ & --- & $\log\eps_{44}$\\
13 & $(-52)\cdot 13$ & --- & $2\log\eps_{13}$\\
15 & $(-15)\cdot 60$ & --- & $4\log\eps_{60}$\\
\bottomrule
\end{tabular}
\end{center}
with $\eps_{12} = 2 + \sqrt3$, $\eps_5 = \tfrac{1+\sqrt5}{2}$, $\eps_{28} =
8 + 3\sqrt7$, $\eps_{44} = 10 + 3\sqrt{11}$, $\eps_{13} =
\tfrac{3+\sqrt{13}}2$, $\eps_{60} = 4 + \sqrt{15}$.  \emph{In every case
$d_2$ is the discriminant of $\Q(\sqrt n)$}: the real quadratic field
carrying the genus-character sum at conductor $n$ is $\Q(\sqrt n)$ itself.
\end{proposition}

\begin{proof}
The identity \eqref{eq:conv} is verified as an identity of integers for
all $m \le 300$ at every level.  This is a proof, by Sturm's bound
\cite{Sturm1987}.  The left side is the $q$-expansion of
$\tfrac12\sum_\idealc\chi_2(\idealc)\theta_{Q_\idealc}$, whose constant
term vanishes and which, by Hecke and Schoeneberg
\cite{Schoeneberg1939}, is a holomorphic modular form of weight $1$ on
$\Gamma_0(|D|)$, $|D| = 4n^2$, with character $\chi_D$; the right side
is the weight-$1$ Eisenstein series $E_1(\chi_{d_1}, \chi_{d_2})$ of
\cite[\S4.8]{DiamondShurman2005}, whose constant term vanishes because
both characters are nontrivial, together with its $p^a$-dilations
weighted by the correction coefficients (at $n = 9$: $E + 3E(9\tau)$,
$E$ of level $36$), a form of the same weight, level and character.  A
weight-$1$ form on $\Gamma_0(4n^2)$ vanishing at $\infty$ to order beyond
$\mu(\Gamma_0(4n^2))/12$ --- at most $180$, attained at $n = 15$, far below
the certified $300$ --- vanishes identically.  For the closed form:
$L(s,\chi_{d_1})$ differs from the primitive $L(s,\chi_{d_1^*})$ by Euler
factors at primes $p \mid d_1$, $p \nmid d_1^*$, which here occur only at
$n = 9$ ($p = 3$, where $\chi_{-3}(3) = 0$ and nothing changes); the
positive $d_2$ is fundamental at every level, $L(0,\chi_{d_2}) = 0$ and
$L'(0,\chi_{d_2}) = h(d_2)\log\eps_{d_2}$, while $L(0,\chi_{d_1^*}) =
2h(d_1^*)/w(d_1^*)$ (Dirichlet's class number formulas at $s = 0$).
\end{proof}

At $n = 9$ the factor $C(s) = 1 + 3^{1-2s}$, $C(0) = 4$, together with
$d_1^* = -3$ ($w = 6$), gives $\tfrac13\cdot 4\log\eps_{12}$.  This
conductor correction at the square conductor --- and only there --- is
thereby proved as an identity; what remains open is its conceptual
derivation from the local proper-$\Oc_9$-ideal structure at $3$, and the
general-$n$ law.  The composite square-free conductor $n = 15$ needs no
correction for its $\Q(\sqrt{15})$-character, as this reading predicts
(verified at $400$ digits).  All seven closed forms agree with the
independent evaluation of Lemma \ref{lem:gamma-eval} to $10^{-249}$ or
better.

\begin{remark}[the other two real characters at $n = 15$; computational]
\label{rem:n15}
The characters $\chi_{(0,1)}$ and $\chi_{(2,1)}$ of $\Cl(\Oc_{15}) \cong
\Z/4\times\Z/2$ admit \emph{no} factorization of the shape \eqref{eq:conv}
with a finite correction over the divisor pairs of $-900$ --- an honest
negative result of the discovery procedure --- but their values are
cleanly identified (certified ratio, $398$ spare digits at $400$ digits):
\[
  L'(0,\chi_{(0,1)}) = 8\log\eps_5, \qquad
  L'(0,\chi_{(2,1)}) = \tfrac83\log\eps_{12},
\]
the units of the two proper subfield levels $5$ and $3$.  Their
factorizations presumably need Euler data at both conductor primes
simultaneously; this is an identification, not a proof.
\end{remark}

\subsection{Euclidean: the $j$-dressing}\label{subsec:dressing-euc}

For the real character $\chi_2$ split the classes by $\chi_2 = \pm1$, and
let $A, B$ be the products of $j(\idealc)$ (resp.\ of $j(\idealc) - 1728$)
over the two cosets, so that $\Sigma_0(\chi_2) = \log|A/B|$ (resp.\
$\Sigma_{1728}(\chi_2) = \log|A/B|$).  By genus theory $A$ and $B$ lie in
the real quadratic subfield $\Q(\sqrt{d_2})$ of the ring class field
fixed by $\ker\chi_2$ and the mirror, are algebraic integers, and are
conjugate over $\Q$, so $A/B$ has norm $1$.

\begin{proposition}[exact $\Sigma$-factorizations; computational]
\label{prop:dressing-euc}
At every level $n \le 15$ and for every real character, $A + B$ and $AB$
are certified integers ($\ge 174$ spare digits at $250$ digits, $\ge 315$
at $400$), $A, B = \tfrac{s \pm t\sqrt{d_2}}2$ in $\Q(\sqrt n)$, and the
norm-one ratio $A/B$ factors \emph{exactly}, by ideal arithmetic in
$\Q(\sqrt{d_2})$ with valuations computed by Hensel lifting and each
identity re-multiplied out symbolically,
\begin{equation}\label{eq:dressing}
  \Sigma_0(\chi_2) = \log\Bigl|\frac AB\Bigr|
  = k\log\eps_{d_2} + \sum_pe_p\,\lambda_p,
  \qquad
  \lambda_p := \log\Bigl|\frac{\pi_p}{\pi_p'}\Bigr|,
\end{equation}
over split primes $p$ of $\Q(\sqrt{d_2})$, $\pi_p$ an explicit generator
of $\idealp_p^{m_p}$ ($m_p$ the order of $\idealp_p$ in the class group,
$e_p \in \tfrac1{m_p}\Z$), normalized so that $0 \le \lambda_p <
2\log\eps_{d_2}$; and likewise for $\Sigma_{1728}$.  The split-prime
support lies in $\GZ(-4n^2,-3)\cup\{p \mid 2n\}$ for $\Sigma_0$ and in
$\GZ(-4n^2,-4)\cup\{p \mid 2n\}$ for $\Sigma_{1728}$, where $\GZ(D, D')$
is the set of primes dividing some $\tfrac{|DD'| - x^2}{4}$
\cite{GrossZagier1985} --- the same Gross--Zagier sets that govern the
denominators of the phase in \PaperI[Prop.~6.26].  The certified
factorizations, in the canonical normalization:
\begin{center}
\footnotesize
\setlength{\tabcolsep}{3pt}
\begin{tabular}{@{}lp{6.2cm}p{4.5cm}@{}}
\toprule
$n$ & $\Sigma_0(\chi_2)$ & $\Sigma_{1728}(\chi_2)$\\
\midrule
3 & $16\log\eps_{12} - 3\lambda_{11} - 3\lambda_{23}$ &
$4\log\eps_{12} + 2\lambda_{11}$\\
5 & $30\log\eps_5 - 3\lambda_{11} + 3\lambda_{59} + 3\lambda_{71}$ &
$28\log\eps_5 + 2\lambda_{11} + 2\lambda_{19}$\\
7 & $24\log\eps_{28} - 3\lambda_{47} - 3\lambda_{83} - 3\lambda_{131}$ &
$18\log\eps_{28} + 2\lambda_3 - 2\lambda_{19}$\\
9 & $16\log\eps_{12} + 3\lambda_{11} - 3\lambda_{23} + 3\lambda_{47} +
3\lambda_{179} - 3\lambda_{227} - 3\lambda_{239}$ &
$24\log\eps_{12} - 2\lambda_{11} - 4\lambda_{23}$\\
11 & $12\log\eps_{44} - 3\lambda_{107} + 3\lambda_{167} + 3\lambda_{263} +
3\lambda_{347} - 3\lambda_{359}$ &
$4\log\eps_{44} + 6\lambda_{19} + 2\lambda_{43}$\\
13 & $18\log\eps_{13} - 3\lambda_{23} + 3\lambda_{107} + 3\lambda_{251} -
3\lambda_{311} + 3\lambda_{443} - 3\lambda_{491} + 3\lambda_{503}$ &
$28\log\eps_{13} - 2\lambda_3 - 4\lambda_{23} - 2\lambda_{43}$\\
15 $(2,0)$ & $42\log\eps_{60} - 3\lambda_{11} + 3\lambda_{59} +
3\lambda_{71} + 3\lambda_{191} + 3\lambda_{419} + 3\lambda_{479} +
3\lambda_{659}$ &
$52\log\eps_{60} + 2\lambda_{11} + 2\lambda_{43} + 2\lambda_{59}$\\
15 $(0,1)$ & $132\log\eps_5 + 3\lambda_{11} + 3\lambda_{59} +
3\lambda_{71} + 3\lambda_{191} - 3\lambda_{419} - 3\lambda_{479} -
3\lambda_{659}$ &
$124\log\eps_5 - 2\lambda_{11} - 4\lambda_{19} + 4\lambda_{31} -
2\lambda_{59}$\\
15 $(2,1)$ & $62\log\eps_{12} - 9\lambda_{11} - 6\lambda_{47} -
3\lambda_{59} - 3\lambda_{71} - 3\lambda_{191} + 3\lambda_{419} +
3\lambda_{479} - 3\lambda_{659}$ &
$16\log\eps_{12} + 6\lambda_{11} + 4\lambda_{23} + 4\lambda_{47} +
2\lambda_{59}$\\
\bottomrule
\end{tabular}
\end{center}
(The generators $\pi_p$ are printed by the script; e.g.\ at $n = 3$,
$\pi_{11} = -4 - \tfrac32\sqrt{12}$, $\pi_{23} = -5 - 2\sqrt{12}$.)
\end{proposition}

The $\pm3$-multiples in $\Sigma_0$ and the even coefficients in
$\Sigma_{1728}$ are the $(4,3)$-exponent shadow of $u^6 =
-\beta^4(\beta - 1728)^3\cdot(\Delta\text{-quotient})$; note that at $n =
15$ all three real characters carry exact $\Sigma$-data although two of
them lack the finite genus factorization of their $L'$.  A general law
for the exponents $k, e_p$ --- a genus-refined Gross--Zagier statement ---
is open (Section \ref{sec:outlook}).

\begin{theorem}[closed form for the genus character]\label{thm:genus-closed}
For the real character $\chi_2$ of $\Cl(\Oc_n)$, $n \le 13$, and for the
$\Q(\sqrt{15})$-character at $n = 15$,
\[
  S(\chi_2) \;=\; -2\,\frac{2h(d_1^*)}{w(d_1^*)}\,h(d_2)\,C(0)\,\log\eps_{d_2}
  \;+\; \tfrac23\,\Sigma_0(\chi_2) + \tfrac12\,\Sigma_{1728}(\chi_2),
\]
with the tabulated exact $\Sigma$-factorizations: an explicit element of
$\Q\log\eps_{\Q(\sqrt n)} + \sum_p\Q\,\lambda_p$ over the Gross--Zagier
split primes.  (The other two real characters at $n = 15$ carry the same
$\Sigma$-closed forms, with the certified $\eps$-ratio of Remark
\ref{rem:n15} standing in for the $L'$-factorization.)
\end{theorem}

\begin{proof}
Theorem \ref{thm:euc-master} with \eqref{eq:Lprime-genus} and
\eqref{eq:dressing}.  Every ingredient is exact; the assembled identity
is additionally checked against the independent evaluation to
$10^{-247}$.
\end{proof}

So the Euclidean phases are an elliptic-unit system \emph{up to an
identified, Gross--Zagier-supported $S$-integer dressing}: the $L'$-part
is a unit statement, and the $j$-part is the price of the singular-moduli
factors in $u^6$.

\begin{remark}[signs]
The factorizations above are of absolute values, and $S(\chi)$ sees only
$|u|$.  The sign of $A/B$ is also computed exactly by the script.  The
sign of the phase itself on the ambiguous classes --- where it is real ---
is a separate matter: on the Euclidean side reality is the center
criterion of \PaperI[Prop.~7.10], and on the hyperbolic side the sign of
$u_f$ on divisor classes is governed by the archimedean law of
\PaperI[Thm.~6.28].  Signs play no further role in this paper.
\end{remark}

\subsection{Hyperbolic genus characters}\label{subsec:genus-hyp}

At $n = 11, 13, 15$ some odd characters of $\Cl(1 - n^2)$ are real.  The
same machinery evaluates everything.

\begin{proposition}[hyperbolic genus factorizations]\label{prop:genus-hyp}
For the odd real characters of the table, the representation numbers
satisfy $R_\chi = 2\,\mathrm{conv}_{d_1,d_2}$ exactly (verified for all $m
\le 300$, proved by the Sturm argument of Proposition
\ref{prop:genus-euc}, the bounds here being $\le 32$; no conductor
correction arises for any odd character), and hence
\[
  L'(0,\chi) = \frac{2h(d_1^*)}{w(d_1^*)}\,h(d_2)\log\eps_{d_2} :
\]
\begin{center}
\small
\begin{tabular}{@{}rllll@{}}
\toprule
$n$ & odd $\chi$ & $d_1\cdot d_2 = 1 - n^2$ & $L'(0,\chi)$ & field\\
\midrule
11 & $\chi_{(0,1)}$ & $(-15)\cdot 8$ & $2\log(1 + \sqrt2)$ & $\Q(\sqrt2)$\\
11 & $\chi_{(1,0)}$ & $(-3)\cdot 40$ & $\tfrac23\log(3 + \sqrt{10})$ & $\Q(\sqrt{10})$\\
13 & $\chi_{(0,1)}$ & $(-7)\cdot 24$ & $\log(5 + 2\sqrt6)$ & $\Q(\sqrt6)$\\
13 & $\chi_{(1,0)}$ & $(-8)\cdot 21$ & $\log\tfrac{5 + \sqrt{21}}2$ & $\Q(\sqrt{21})$\\
15 & $\chi_{(0,1)}$ & $(-8)\cdot 28$ & $\log(8 + 3\sqrt7)$ & $\Q(\sqrt7)$\\
15 & $\chi_{(2,1)}$ & $(-4)\cdot 56$ & $\tfrac12\log(15 + 4\sqrt{14})$ & $\Q(\sqrt{14})$\\
\bottomrule
\end{tabular}
\end{center}
\end{proposition}

The \emph{even} real characters pair with the complementary fields.  At
$n = 11$ and $n = 13$ the even characters factor (certified by the same
procedure) as
\[
  L(s,\chi_{-24})\,L(s,\chi_5)\quad (n = 11),
  \qquad
  L(s,\chi_{-3})\,L(s,\chi_{56})\quad (n = 13),
\]
so the fields are $\Q(\sqrt5)$ and $\Q(\sqrt{14})$ --- exactly the fields
of the pair-sums $u_f + u_f^{-1}$ of \PaperI[Thm.~6.19(3)].  So the
parity $\chi(\idealr_n) = \pm1$ splits the genus fields of $1 - n^2$ into
the \emph{pair-sum fields} (even; the algebra of $u + 1/u$) and the
\emph{$\log|u|$-fields} (odd; the character sums) --- complementary
halves of one genus theory.  At $n = 9$ (disc $-80$, conductor $2$) the
single real character is even, and at $n = 15$ the even character
$\chi_{(2,0)}$ appears; their Epstein $L$-functions factor as
\begin{gather*}
  L(s,\chi_{-16})\,L(s,\chi_5)\,(1 + 2^{-s} + 2\cdot4^{-s})\quad (n = 9),\\
  L(s,\chi_{-7})\,L(s,\chi_{32})\,(1 - 2^{-s} + 2\cdot4^{-s})\quad (n = 15),
\end{gather*}
two more instances of the conductor Euler phenomenon, both certified.

\begin{proposition}[hyperbolic $\Sigma$-factorizations; computational]
\label{prop:dressing-hyp}
For the odd real characters the coset products over $\ker\chi$ are
certified conjugate quadratic integers in the matching field, and the
ratios factor exactly as in \eqref{eq:dressing}; where a split prime is
non-principal the generator is taken for $\idealp^m$ ($m$ the class
order) and the exponent is $e/m$ --- at $n = 11$ in $\Q(\sqrt{10})$ (class
number $2$), $\pi_3 = 1 + \tfrac12\sqrt{40}$ generates $\idealp_3^2$.
\begin{center}
\footnotesize
\setlength{\tabcolsep}{3pt}
\begin{tabular}{@{}llp{4.8cm}p{5.3cm}@{}}
\toprule
$n$ & $\chi$ & $\Sigma_0(\chi)$ & $\Sigma_{1728}(\chi)$\\
\midrule
11 & $(-15)\cdot8$ & $30\log\eps_8 + 3\lambda_{41} + 3\lambda_{89}$ &
$36\log\eps_8 + 2\lambda_{71}$\\
11 & $(-3)\cdot40$ & $14\log\eps_{40} + 3\lambda_3 - 3\lambda_{41} -
3\lambda_{89}$ & $16\log\eps_{40} - 2\lambda_{71}$\\
13 & $(-7)\cdot24$ & $18\log\eps_{24} - 3\lambda_5 + 3\lambda_{101}$ &
$12\log\eps_{24} + 4\lambda_{19} + 2\lambda_{47} + 2\lambda_{167}$\\
13 & $(-8)\cdot21$ & $24\log\eps_{21} - 3\lambda_5 - 3\lambda_{101}$ &
$24\log\eps_{21} - 2\lambda_{47} - 2\lambda_{167}$\\
15 & $(-8)\cdot28$ & $30\log\eps_{28} - 6\lambda_{29} + 3\lambda_{47} -
6\lambda_{53} + 3\lambda_{167}$ & $20\log\eps_{28} + 2\lambda_{31} -
2\lambda_{47} - 2\lambda_{103} - 2\lambda_{199} - 2\lambda_{223}$\\
15 & $(-4)\cdot56$ & $6\log\eps_{56} - 6\lambda_{11} + 3\lambda_{47} -
3\lambda_{167}$ & $2\log\eps_{56} + 4\lambda_{11} + 2\lambda_{31} +
4\lambda_{43} - 2\lambda_{47} + 2\lambda_{103} - 2\lambda_{199} -
2\lambda_{223}$\\
\bottomrule
\end{tabular}
\end{center}
The split-prime supports again land in $\GZ(1 - n^2, -3)$ resp.\
$\GZ(1 - n^2, -4)$ together with $p \mid 2(n^2 - 1)$.  Combining with
Proposition \ref{prop:genus-hyp} and Theorem \ref{thm:hyp-master},
$S(\chi)$ is in closed form for every odd real character of every
computed level, verified to $10^{-247}$.
\end{proposition}

\emph{The non-real odd characters.}  At $n = 9$ ($\Cl \cong \Z/4$) the
two odd characters are quartic, and their sums collapse by the twist law
to $S(\chi_4) = 2\log|u_1|$: the master identity is then precisely
Lemma \ref{lem:hyp-classwise} at the principal class, and $-12L'(0,\chi_4)
= \log|R_1|$ is a Stark-type evaluation in the quartic ring class field,
certified through the polynomial of Theorem \ref{thm:Rpoly}.  The same
holds at $n = 15$ for its quartic odd pair.

\subsection{Certified non-fits: the Stark regime}\label{subsec:nonfits}

\begin{proposition}[computational; certified negative results]
\label{prop:nonfits}
For the order-$3$ and order-$6$ characters at Euclidean $n = 9, 11, 13$,
the order-$4$ characters at Euclidean $n = 7, 15$, and the odd quartic
characters at hyperbolic $n = 9, 15$, an integer-relation search in the
information-theoretically safe regime ($250$ digits, $400$ at Euclidean
$n = 15$; coefficient height $\le 10^6$; tolerance $10^{-200}$, resp.\
$10^{-320}$; a relation is accepted only when its information content is
at most a quarter of the working precision) finds \emph{no} rational
linear relation among $S(\chi)$, $L'(0,\chi)$, $\{\log p : p \mid 2nH(0)
H(1728)\}$ (resp.\ $p \mid 2(n^2-1)H(0)H(1728)$) and $\log\eps_{d_2}$.
\end{proposition}

The dressing of a non-real character genuinely lives in the cubic or
quartic subfield of the ring class field --- its exact description is the
certified coset data of Remark \ref{rem:cosetcubics} --- and not in the
rank-one basis of the naive conjecture.  This is the expected Stark-regime
behaviour: the master identities are exactly the part that survives, and
the recognition of the coset objects as Stark units of the cubic and
quartic subfields \cite[Ch.~6--7]{Schertz2010} is the natural
continuation (Section \ref{sec:outlook}).

%%% ============================================================
\section{The unit theorem and the per-class valuation law}\label{sec:unit}

This is the technical core of the paper.  Throughout, $n \ge 3$ is odd,
$D = 1 - n^2$, $\Oc$ the order of discriminant $D$, and the notation of
the introduction is in force: $r_0 = \tfrac{n-1}2$, $\idealr = \lat{r_0,
\omega_0}$, $\idealb_1 = \lat{1, m_1}$, $\beta_1 = j(f)$, $\beta_2 =
j(\idealr f)$, and $R_f$ as in \eqref{eq:Rdef}.

\subsection{The lattice form of the twisted $\Delta$-ratio}
\label{subsec:lattice}

\begin{lemma}[lattice lemma]\label{lem:lattice}
For every primitive class $f$ of every odd level $n \ge 3$,
\begin{equation}\label{eq:lattice}
  R_f \;=\; r_0^{\,6}\;\frac{\Delta(\idealb_1)}{\Delta(\idealr^{-1}\idealb_1)} .
\end{equation}
In particular the $j$-dressing of $R_f$ is pure bookkeeping: the twisted
ratio \emph{is} the Siegel-type $\Delta$-quotient along the invertible
ambiguous twist, normalized by its norm.
\end{lemma}

\begin{proof}
From the closed form $u_f = -\eps\mu^{-2}h_2(\idealb_1)/h_2(\idealb_2)$
\PaperI[Prop.~6.10] and the weight algebra $h_2^6 = c\,j^4(j-1728)^3
\Delta$,
\[
  R_f = \eps^6\mu^{-12}\,\frac{\Delta(\idealb_1)}{\Delta(\idealb_2)} .
\]
Since $[\idealb_2] = [\idealr^{-1}\idealb_1]$ there is $\nu \in K^\times$,
unique up to sign, with $\idealb_2 = \nu\,\idealr^{-1}\idealb_1$
\PaperI[eq.~(6.1)]; homogeneity gives
\[
  \Delta(\idealb_2) = \nu^{-12}\Delta(\idealr^{-1}\idealb_1),
  \qquad\text{so}\qquad
  R_f = \eps^6\Bigl(\frac{\nu^2}{\mu^2}\Bigr)^{6}\,
  \frac{\Delta(\idealb_1)}{\Delta(\idealr^{-1}\idealb_1)} .
\]
By the theorem
$\omega_f = \eps\nu^2/(r_0\mu^2) = 1$ of \PaperI[Thm.~6.18],
$\eps\nu^2/\mu^2 = r_0$, whence $\eps^6(\nu^2/\mu^2)^6 = r_0^6$.
\end{proof}

The lemma is verified classwise at every odd $n \le 21$, with maximal
relative deviation $6.5\cdot 10^{-419}$ at $420$ digits, the twist
lattice $\idealr^{-1}\idealb_1$ being computed by exact Hermite-normal-form
arithmetic in $K$.  At $n = 3$: $r_0 = 1$, $\idealr = \Oc$, and $R = 1$
exactly --- the mandatory anchor.

\subsection{The unit theorem}\label{subsec:unit}

\begin{theorem}[unit theorem]\label{thm:unit}
For every odd $n \ge 3$ and every class $f$ of discriminant $1 - n^2$ ---
primitive or not, with $R_f$ defined for imprimitive classes by the
lattice formula \eqref{eq:lattice} --- one has $v_\idealP(R_f) = 0$ at
every finite place $\idealP$ of $\bar\Q$: $R_f$ is an algebraic
\emph{unit}.  Consequently the level polynomial $\prod_f(x - R_f)$ over
the primitive classes is a monic integer polynomial with constant term
$\pm1$, palindromic by the twist law: the certified observations of
Theorem \ref{thm:Rpoly} are theorems for every $n$.
\end{theorem}

The proof runs through five lemmas.  Fix a prime $p$ and a place
$\idealP \mid p$ of $\bar\Q$, with $v$ the valuation normalized by
$v(p) = 1$.  For a lattice $\ideala$ with CM by $\Oc$ write $E_\ideala$
for the elliptic curve $\C/\ideala$.

\begin{lemma}[tangent-scalar expression]\label{lem:tangent}
Let $\ideala$ be a lattice with CM by $\Oc$, $\idealc$ an invertible
$\Oc$-ideal, and $\varphi: E_\ideala \to E_{\idealc^{-1}\ideala}$ the
isogeny that is $z \mapsto z$ analytically.  Then
\[
  v\Bigl(\frac{\Delta(\ideala)}{\Delta(\idealc^{-1}\ideala)}\Bigr)
  \;=\; -12\,v(\delta_\varphi),
\]
where $\delta_\varphi$ is the scalar by which $\varphi$ pulls back N\'eron
differentials of good-reduction models at $\idealP$.
\end{lemma}

\begin{proof}
Choose Weierstrass models $(E_i, \omega_i)$ over $\bar\Q$ whose period
lattices are $\lambda\ideala$ and $\mu\,\idealc^{-1}\ideala$; then
$\Delta(E_1, \omega_1) = \lambda^{-12}\Delta(\ideala)$, $\Delta(E_2,
\omega_2) = \mu^{-12}\Delta(\idealc^{-1}\ideala)$, and $\varphi^*\omega_2
= t\,\omega_1$ with $t = \mu/\lambda \in \bar\Q$ (the isogeny and the
differentials are algebraic), so $\Delta(\ideala)/\Delta(\idealc^{-1}
\ideala) = t^{-12}\Delta(E_1, \omega_1)/\Delta(E_2, \omega_2)$.  CM curves
have potentially good reduction; over a finite extension $L/\Q_p$ choose
minimal models: their discriminants are units, and writing $\omega_i =
u_i\,\omega_{i,\min}$ one gets $v(\Delta(E_i, \omega_i)) = -12v(u_i)$
and $\delta_\varphi = t\,u_1/u_2$; the $u_i$ cancel and the display
follows.
\end{proof}

\begin{lemma}[multiplicativity; prime-to-$p$ triviality]\label{lem:mult}
Tangent scalars multiply along composites, and $v(\delta_\varphi) = 0$
whenever $\deg\varphi$ is prime to $p$.
\end{lemma}

\begin{proof}
Multiplicativity is the chain rule for pullbacks.  Both $\delta_\varphi$
and $\delta_{\hat\varphi}$ are integral --- isogenies extend to N\'eron
models --- and $\delta_\varphi\delta_{\hat\varphi} = \deg\varphi$, a
$p$-adic unit.
\end{proof}

\begin{lemma}[the $p$-part of the twist]\label{lem:ppart}
$\idealr = \idealb\idealc$ with $\idealb, \idealc$ invertible
$\Oc$-ideals, $N\idealb = p^k$ where $k = v_p(r_0)$, $N\idealc$ prime to
$p$, and $\idealb$ inherits both structural properties of $\idealr$:
$\bar\idealb = \idealb$ and $\idealb^2 = (p^k)$.
\end{lemma}

\begin{proof}
Invertible fractional $\Oc$-ideals decompose along primes (localization);
let $\idealb$ be the invertible ideal with $\idealb_q = \Oc_q$ for $q \ne
p$ and $\idealb_p = \idealr_p$, and $\idealc = \idealb^{-1}\idealr$.
Conjugation acts locally, and $\bar\idealr = \idealr$ gives $\bar\idealb
= \idealb$.  Locally $(\idealb^2)_p = (\idealr^2)_p = r_0\Oc_p =
p^k\Oc_p$ and $(\idealb^2)_q = \Oc_q$, which is the local data of the
principal ideal $(p^k)$.
\end{proof}

\begin{lemma}[rigidity]\label{lem:rigidity}
The quantity $d_p := v(\delta_\varphi)$ for the $\idealb$-twist isogeny
$\varphi: E_\ideala \to E_{\idealb^{-1}\ideala}$ is the same for every
proper $\Oc$-ideal $\ideala$.
\end{lemma}

\begin{proof}
$v(\delta_\varphi)$ is computed on formal groups: it is the tangent
valuation of the quotient map $G \to G/C$ of $p$-divisible groups over
$\Oc_L$, where $G = E_\ideala[p^\infty]$ (of the good model) and $C$ is the
schematic closure of $\ker\varphi = \idealb_p^{-1}\ideala_p/\ideala_p$ ---
the quotient elliptic curve is the good model of $E_{\idealb^{-1}\ideala}$
by the N\'eron property, and $\Lie$ of an elliptic curve is $\Lie$ of its
$p$-divisible group.  So $d_p$ depends only on the isomorphism class of
the pair $(G, C)$.  Given two proper ideals, choose an invertible
$\Oc$-ideal $\idealq$ of norm prime to $p$ connecting their classes; the
connecting isogeny has degree prime to $p$, hence induces an isomorphism
of $p$-divisible groups, is $\Oc_p$-linear on Tate modules ($T_pE_\ideala
= \ideala\otimes\Z_p$, free of rank one over $\Oc_p$ since proper ideals
of quadratic orders are invertible, hence locally principal), and
identifies the two kernels $\idealb_p^{-1}\ideala_p/\ideala_p$.
Homothety adjusts representatives within a class.
\end{proof}

\begin{proposition}[balance]\label{prop:balance}
$2\,d_p = k$.
\end{proposition}

\begin{proof}
Compose the $\idealb$-twist at $\ideala$ with the $\idealb$-twist at
$\idealb^{-1}\ideala$: since $\idealb^2 = (p^k)$, $\idealb^{-2}\ideala =
p^{-k}\ideala$, and the composite followed by the isomorphism $z \mapsto
p^kz$ is the endomorphism $[p^k]$ of $E_\ideala$, of tangent scalar
$p^k$.  By Lemma \ref{lem:mult} the tangent valuations add to $k$
(an isomorphism of good models contributes $0$); by Lemma
\ref{lem:rigidity} both steps contribute the same $d_p$, the twist of a
proper ideal being proper.
\end{proof}

\begin{proof}[Proof of Theorem \ref{thm:unit}]
By Lemma \ref{lem:lattice} and Lemma \ref{lem:tangent} applied to the
$\idealr$-twist $\varphi: E_{\idealb_1} \to E_{\idealr^{-1}\idealb_1}$,
which by Lemma \ref{lem:ppart} is the composite of the $\idealb$-twist
and a twist of degree $N\idealc$ prime to $p$,
\[
  v(R_f) = 6\,v(r_0) - 12\,v(\delta_\varphi) = 6k - 12\,d_p = 6k - 6k = 0
\]
at every $\idealP \mid p$ with $p \mid r_0$ (Lemmas \ref{lem:mult},
\ref{lem:rigidity} and Proposition \ref{prop:balance}), and $v(R_f) = 0$
trivially at $p \nmid r_0$, where the whole twist has degree prime to
$p$.  For an imprimitive class of content $g$, $\idealb_1$ is proper over
the order $\Oc'$ of discriminant $D/g^2$ and the twist is by $\idealr' :=
\idealr\Oc'$ --- invertible (the extension of an invertible ideal),
ambiguous, with $\idealr'^2 = r_0\Oc'$ --- so the identical argument runs
over $\Oc'$.  Finally the polynomial statement: rationality of
$\prod_f(x - R_f)$ is Theorem \ref{thm:Rpoly}; the roots being units, the
coefficients are rational algebraic integers, hence integers, with
$\prod_fR_f = 1$ for $n \ge 5$ (the twist law pairs $f$ with $\idealr f$)
and $R = 1$ at $n = 3$; palindromy is the twist law $R_{\idealr f} =
R_f^{-1}$.
\end{proof}

\begin{remark}[what the proof did not need]\label{rem:notneeded}
No Kronecker limit formula, no Siegel units, and --- despite the Deuring
dictionary that motivated the statement --- no quasi-canonical-lifting
valuation formulas: Gross's theory \cite{Gross1986} enters only as the
conceptual reading (the twist by an ambiguous ideal of the order is a
\emph{horizontal} move, and horizontal moves preserve local
$\Delta$-data).  The inputs are potentially good reduction, N\'eron
differentials, and the rigidity of $p$-divisible-group quotients.
\end{remark}

\subsection{The certified record and the imprimitive strata}
\label{subsec:unitrecord}

\begin{proposition}[level polynomials; computational]\label{prop:Rrecord}
For every odd $n \le 21$ the level polynomial $\prod_{f\ \mathrm{prim}}
(x - R_f)$ is certified with integer coefficients, palindromic, with
constant term $1$ (spare digits $\ge 237$ at $250$ digits for $n \le 13$,
$\ge 275$ at $300$ digits for $n = 15, 17, 19$, $385$ at $420$ digits for
$n = 21$) --- as Theorem \ref{thm:unit} guarantees --- and is moreover
\emph{irreducible over $\Q$} (exact factorization): the $R_f$ of a level
form a single Galois orbit.
\begin{center}
\footnotesize
\setlength{\tabcolsep}{4pt}
\begin{tabular}{@{}rrp{11.1cm}@{}}
\toprule
$n$ & $h$ & $\prod_f(x - R_f)$\\
\midrule
3 & 1 & $x - 1$\\
5 & 2 & $x^2 - 34x + 1$\\
7 & 2 & $x^2 - 2702x + 1$\\
9 & 4 & $x^4 - 339524x^3 - 95354x^2 - 339524x + 1$\\
11 & 4 & $x^4 - 56529284x^3 + 1538876166x^2 - 56529284x + 1$\\
13 & 4 & $x^4 - 11382984004x^3 + 885435408006x^2 - 11382984004x + 1$\\
15 & 8 & $x^8 - 2628641876392x^7 - 21595933374628x^6 -
1373071731101336x^5 + 9740462908109254x^4 - 1373071731101336x^3 -
21595933374628x^2 - 2628641876392x + 1$\\
17 & 4 & $x^4 - 673122277718404x^3 + 8553847041196806x^2 -
673122277718404x + 1$\\
19 & 8 & $x^8 - 186863535844922888x^7 + 44665915402536486036508x^6 +
131633377547326082495944x^5 + 179879784238619113420870x^4 +
(\text{sym})$\\
21 & 12 & $x^{12} - 55348592922774901452x^{11} +
87282798056992201360611266x^{10} + \cdots$ (all coefficients printed by
the script)\\
\bottomrule
\end{tabular}
\end{center}
\end{proposition}

Irreducibility for general $n$ remains open exactly as for the phase
polynomials $Q_n$ of \PaperI[Open problem~6.24]: the coincidence set
$\{f : R_f = R_1\}$ is a translation-invariant subgroup, and only its
triviality is missing.

\begin{proposition}[imprimitive strata]\label{prop:strata}
Let $f$ be a class of content $g > 1$, $\Oc'$ the order of discriminant
$D/g^2$, and $\idealr' = \idealr\Oc'$ the induced twist ideal.  With $R$
defined by the lattice formula, every stratum polynomial is again an
integer palindromic polynomial with constant term $\pm1$ (Theorem
\ref{thm:unit}), and
\[
  \idealr'\ \text{principal in }\Oc' \;\Longrightarrow\; R \equiv 1
  \text{ on the content-}g\text{ stratum}.
\]
The converse holds on every computed stratum (computational).
\end{proposition}

\begin{proof}
If $\idealr' = (\lambda)$, then $\lambda^2\Oc' = \idealr'^2 = r_0\Oc'$
forces $\lambda^2 = r_0\zeta$ with $\zeta \in \Oc'^\times$ a root of
unity, so $\Delta(\idealr'^{-1}\idealb_1) = \lambda^{12}\Delta(\idealb_1)
= r_0^6\Delta(\idealb_1)$ and $R = 1$.
\end{proof}

\begin{center}
\footnotesize
\setlength{\tabcolsep}{3pt}
\begin{tabular}{@{}rrrll@{}}
\toprule
$n$ & content & disc & $\idealr\Oc'$ & stratum polynomial\\
\midrule
7 & 2 & $-12$ & $(\sqrt{-3})$, principal & $x - 1$\\
7 & 4 & $-3$ & principal & $x - 1$\\
9 & 2 & $-20$ & $(2)$, principal & $(x - 1)^2$\\
15 & 2 & $-56$ & $\lat{7, \sqrt{-14}}$, \textbf{non}principal &
$x^4 - 1988x^3 - 3194x^2 - 1988x + 1$\\
17 & 2 & $-72$ & $\lat{4, 6\sqrt{-2}}$, \textbf{non}principal &
$x^2 - 9602x + 1$\\
17 & 3 & $-32$ & $(2\sqrt{-2})$, principal & $(x - 1)^2$\\
17 & 6 & $-8$ & $(2\sqrt{-2})$, principal & $x - 1$\\
19 & 3 & $-40$ & $(3)$, principal & $(x - 1)^2$\\
\bottomrule
\end{tabular}
\end{center}
So nothing breaks on the imprimitive strata: the unit theorem persists;
what changes is that the twist can \emph{die} (principal induced ideal),
collapsing the stratum to $R \equiv 1$.

\subsection{The per-class valuation law}\label{subsec:perclass}

We return to the Euclidean $\Delta$-data $G_\idealc = n^{12}\Delta(
\Lambda_\idealc)/\Delta(\Zi)$ over $\Cl(\Oc_n)$, with $D_n(x) =
\prod_\idealc(x - G_\idealc) \in \Z[x]$ and $|D_n(0)| = |M(n)|$ (Theorem
\ref{thm:Dn}).  The $\Delta$-mass law \PaperI[Thm.~7.18] controls the
product of the $G_\idealc$; the rigidity of Lemma \ref{lem:rigidity}
distributes it over the classes.

\begin{theorem}[per-class valuations]\label{thm:perclass}
For every level $n \ge 2$, every prime $p$ with $p^k \parallel n$, and
every place $\idealP \mid p$ of $\bar\Q$ (normalized $v(p) = 1$), the
valuation $v_\idealP(G_\idealc)$ is \emph{independent of the class}
$\idealc$, equal to
\begin{equation}\label{eq:wpk}
  w_p(k) \;=\;
  \begin{cases}
  0, & p \text{ split } (p \equiv 1 \bmod 4),\\[2pt]
  \dfrac{12\,(p^k - 1)}{(p-1)\,p^{k-1}(p+1)}, & p \text{ inert }
  (p \equiv 3 \bmod 4),\\[6pt]
  6\,(2^k - 1)/2^k, & p = 2;
  \end{cases}
\end{equation}
and $v_\idealP(G_\idealc) = 0$ at $p \nmid n$.
\end{theorem}

\begin{proof}
Write $\psi_\idealc: E_{\Lambda_\idealc} \to E_0 := E_{\Zi}$ for the
isogeny $z \mapsto z$ (degree $n$); by Lemma \ref{lem:tangent},
$v(G_\idealc) = 12k - 12\,v(\delta)$, where $\delta$ is the tangent scalar
of the $p$-primary part of $\psi_\idealc$ (Lemma \ref{lem:mult}).  By
duality it suffices to control $\hat\psi$-data on the \emph{fixed} curve
$E_0$: the $p$-primary kernel of the dual is the subgroup $C_\idealc =
p^{-k}\Lambda_{\idealc,p}/\Z_p[i] \subset E_0[p^\infty]$.  Local structure:
$\Lambda_{\idealc,p}$ is an invertible $\Oc_{n,p}$-module inside $\Z_p[i]$
of index $p^k = [\Z_p[i] : \Oc_{n,p}]$, hence $\Lambda_{\idealc,p} =
\epsilon\,\Oc_{n,p}$ with $\epsilon \in K_p^\times$, and the containment
and the index force $\epsilon \in \Z_p[i]^\times$.  Since $\Z_p[i]
\subset \operatorname{End}(E_0[p^\infty])$ (the closure of $\Zi$ in the
$p$-adically complete endomorphism ring), $\epsilon$ is an
\emph{automorphism of the $p$-divisible group} carrying $C_{\idealc'}$ to
$C_\idealc$: the pairs $(E_0[p^\infty], C_\idealc)$ are isomorphic for all
classes, so $v(\delta_{\hat\psi_\idealc})$, hence $v(G_\idealc)$, is
class-independent.  The value is then pinned by the $\Delta$-mass law:
$h\cdot w = v_p(\prod_\idealc G_\idealc) = v_p(M(n))$, and $v_p(M(n))/h =
w_p(k)$ by \PaperI[Thm.~7.18] together with $h = \tfrac12\Ne(p^k)
\Ne(n/p^k)$ \PaperI[Prop.~3.1, Thm.~7.6].
\end{proof}

\begin{proposition}[Newton-polygon verification; computational]
\label{prop:newton}
The multiset of $\idealP$-adic valuations of the roots of $D_n$ is its
Newton polygon at $p$.  For every computed level --- $2 \le n \le 18$,
$21$, $25$, $27$, $49$ --- and every $p \mid n$, the polygon of the
certified $D_n$ (spare digits $140$--$249$; $460$ digits at $n = 49$) has a
\emph{single slope, of the value \eqref{eq:wpk}}, across the full degree
$h$:
\begin{center}
\small
\begin{tabular}{@{}rl@{\qquad}rl@{}}
\toprule
$n$ & polygon slopes & $n$ & polygon slopes\\
\midrule
2 & $2 : 3$ & 12 & $2 : 9/2;\ 3 : 3$\\
4 & $2 : 9/2$ & 14 & $2 : 3;\ 7 : 3/2$\\
7 & $7 : 3/2$ & 16 & $2 : 45/8$\\
8 & $2 : 21/4$ & 18 & $2 : 3;\ 3^2 : 4$\\
9 & $3^2 : 4$ & 27 & $3^3 : 13/3$\\
11 & $11 : 1$ & 49 & $7^2 : 12/7$\\
\bottomrule
\end{tabular}
\end{center}
(split primes: slope $0$ throughout, e.g.\ $n = 5, 13, 25$; at composite
$n$ every prime factor is listed).
\end{proposition}

The denominators of the slopes --- $p^{k-1}(p+1)$ at inert $p$, $2^k$ at
the ramified prime --- are exactly the ramification degrees of Gross's
quasi-canonical liftings of level $k$ \cite{Gross1986}: the per-class law
\emph{is} the quasi-canonical valuation ladder, obtained here from
rigidity plus the elementary mass recursion.  We stress that this reading
is an interpretation and not an input of the proof.

\begin{corollary}[the split ladder]\label{cor:split}
Let $p \equiv 1 \pmod 4$, $p \nmid 6$, $p^k \parallel n$, $n \ne p^k$.
Then for \emph{every} class, $v_\idealP(u_\idealc^2) = -4k$.
\end{corollary}

\begin{proof}
By \PaperI[Lemma~7.12],
\[
  u^2 = -12\,\beta(\beta - 1728)\,\frac{g_2(\Lambda_\idealc)}{g_2(\Zi)},
  \qquad
  \Bigl(\frac{g_2(\Lambda_\idealc)}{g_2(\Zi)}\Bigr)^{3}
  = \frac{\beta}{1728}\cdot\frac{G_\idealc}{n^{12}},
\]
so $v(u^2) = \tfrac43v(\beta) + v(\beta - 1728) + \tfrac13v(G_\idealc) -
4k$.  Theorem \ref{thm:perclass} gives $v(G_\idealc)
= 0$.  At an ordinary place, $\beta \equiv 0$ would force the
endomorphism ring of the reduction to contain orders of two distinct
imaginary quadratic fields (impossible), and $\beta \equiv 1728$ would
force equal prime-to-$p$ conductors, $n/p^k = 1$ (Deuring), which is
excluded.
\end{proof}

This proves the exactly arithmetic $5^{4k}$-denominator ladder observed
at $n = 15$ in \PaperI[\S7.4] (with the witness $v_5(H_{-900}(0)) =
v_5(H_{-900}(1728)) = 0$ certified).  At non-split $p$ the same reduction
expresses the fine structure of the normalizers $\lambda_n$ of
\PaperI[\S7.4] through $w_p(k)$ plus the $\beta$-side collision
valuations; the remaining open ingredient is the conductor-degenerate
Gross--Zagier data of Section \ref{sec:outlook}, no longer the
$\Delta$-part.

\begin{corollary}[the inert ladders, decomposed; computational]
\label{cor:inert}
At $n = 7$ ($p = 7$, $k = 1$) and $n = 9$ ($p = 3$, $k = 2$) every
ingredient of $v_\idealP(u_\idealc^2) = \tfrac43v(\beta) + v(\beta - 1728) +
\tfrac13v(G) - 4k$ is certified to be a \emph{single} Newton-polygon
slope --- the $\beta$-collision valuations are class-independent too ---
giving the constant per-class values
\begin{gather*}
  n = 7:\quad v_7(u_\idealc^2) = \tfrac43\cdot0 + \tfrac14 + \tfrac13\cdot\tfrac32 - 4
  = -\tfrac{13}{4},\\
  n = 9:\quad v_3(u_\idealc^2) = \tfrac43\cdot\tfrac12 + \tfrac12 + \tfrac13\cdot4 - 8
  = -\tfrac{11}{2}.
\end{gather*}
The recorded denominator ladders follow: at $n = 7$ the integrality bound
$7^{\lfloor 13k/4\rfloor} = 7^3, 7^6, 7^9, 7^{13}$ for the elementary
symmetric functions $e_k$ of the $u_\idealc^2$ is met \emph{exactly}; at
$n = 9$ the bound $3^{\lfloor 11k/2\rfloor}$ is attained at the
coset-structured $k = 3, 6$ ($3^{16}, 3^{33}$), with structured
cancellation at the other $k$ ($3^2, 3^6, 3^{19}, 3^{22}$).
\end{corollary}

So the $\lambda_n$-question is reduced, at every computed level, to one
datum: the per-class $\beta$-collision slopes, certified constant here,
whose general law is the genus-refined Gross--Zagier problem of Section
\ref{sec:outlook}.

%%% ============================================================
\section{First-power Schmidt units}\label{sec:firstpower}

The unit $R_f$ is a sixth-power-level object.  A canonical sixth root
exists and obeys the laws of the phase at first power.  Let $\gamma_2 =
E_4/\eta^8$ and $\gamma_3 = E_6/\eta^{12}$ be the canonical holomorphic
Weber functions ($\gamma_2^3 = j$, $\gamma_3^2 = j - 1728$), and for a
primitive class $f$ let $\tau_f$ be the CM point of its reduced form.

\begin{definition}\label{def:w}
$\displaystyle
  w_f \;:=\; u_f\cdot\frac{(\gamma_2^2\gamma_3)(\tau_{\idealr f})}
  {(\gamma_2^2\gamma_3)(\tau_f)} .$
\end{definition}

\begin{proposition}[exact first-power laws]\label{prop:wlaws}
For every odd $n \ge 3$ and every primitive class $f$:
\begin{enumerate}
\item $w_f^{\,6} = R_f$;
\item $w_f\,w_{\idealr f} = 1$;
\item $\overline{w_f} = w_{f^{-1}}$ whenever neither reduced
representative is a boundary form ($b = a$ or $a = c$);
\item $\sigma(w_f) = \zeta(\sigma, f)\,w_{f^{e(\sigma)}\idealc(\sigma)}$
with $\zeta(\sigma, f)^6 = 1$, for every $\sigma \in \Gal(\bar\Q/\Q)$.
\end{enumerate}
Moreover, let $\tau_0$ be the adapted uniformizer of \PaperI[Lemma~6.13]
for the pair $(\idealb_1, \idealr^{-1}\idealb_1)$; by \PaperI[Thm.~6.18],
\[
  u_f = -\,\frac{1}{r_0}\,\frac{j'(\tau_0)}{j'(r_0\tau_0)} .
\]
Then
\begin{equation}\label{eq:etaquot}
  w_f \;=\; -\,\frac{1}{r_0}\Bigl(\frac{\eta(\tau_0)}{\eta(r_0\tau_0)}\Bigr)^{\!4}
  \quad\text{up to a sixth root of unity}:
\end{equation}
a weight-$0$ $\eta$-quotient on $\Gamma_0(r_0)$ evaluated at a Heegner
point, well defined up to the $\mu_6$-valued $\eta^4$-multiplier of the
basis choice.  The cocycle $\zeta(\sigma, f)$ \emph{is} this multiplier
system.
\end{proposition}

\begin{proof}
(1) $(\gamma_2^2\gamma_3)^6 = j^4(j - 1728)^3$, so $w_f^6 = u_f^6\,
\beta_2^4(\beta_2 - 1728)^3/\beta_1^4(\beta_1 - 1728)^3 = R_f$.
(2) The twist law $u_fu_{\idealr f} = 1$ \PaperI[Thm.~6.6] and the
telescoping of the $\gamma$-ratio along the $2$-torsion twist
($\idealr^2f = f$, and reduced points are canonical).
(3) Off the boundary of the reduced domain $\tau_{f^{-1}} = -\bar\tau_f$
exactly, the $q$-expansions of $\gamma_2, \gamma_3$ are real, so
$(\gamma_2^2\gamma_3)(-\bar\tau) = \overline{(\gamma_2^2\gamma_3)(\tau)}$,
and the mirror law $u_{f^{-1}} = \bar u_f$ \PaperI[Thm.~6.6] finishes.
(4) Sixth powers are strictly equivariant by Theorem \ref{thm:unit}
together with the first-power descent \PaperI[Thm.~6.19] (the argument
of Theorem \ref{thm:Rpoly}), so the discrepancy at first power is a
sixth root of unity.  For \eqref{eq:etaquot}: $j' = -2\pi i\,
\gamma_2^2\gamma_3\,\eta^4$, and $\tau_f, \tau_{\idealr f}$ are
$\SL_2(\Z)$-equivalent to $\tau_0, r_0\tau_0$, under which
$\gamma_2^2\gamma_3$ picks up sixth roots of unity.
\end{proof}

\begin{proposition}[the coherence table; computational]\label{prop:mn}
Let $m(n) \mid 6$ be minimal with $\{w_f^{\,m}\}$ Galois-stable
(equivalently, $\prod_f(x - w_f^{\,m})$ a certified integer polynomial).
For all odd $3 \le n \le 35$:
\begin{center}
\footnotesize
\setlength{\tabcolsep}{3.2pt}
\begin{tabular}{@{}l*{17}{c}@{}}
\toprule
$n$ & 3 & 5 & 7 & 9 & 11 & 13 & 15 & 17 & 19 & 21 & 23 & 25 & 27 & 29 &
31 & 33 & 35\\
\midrule
$m(n)$ & 1 & 3 & 1 & 3 & 6 & 1 & 1 & 3 & 2 & 3 & 6 & 1 & 6 & 6 & 6 & 6 & 6\\
\bottomrule
\end{tabular}
\end{center}
At the fully coherent levels the first-power Schmidt units exist
unconditionally, with certified irreducible minimal polynomials
\begin{align*}
n = 3:&\quad x + 1 \qquad (w = -1),\\
n = 7:&\quad x^2 + 4x + 1 \qquad (w = -\eps_{12}^{\mp1} = -2 \pm \sqrt3),\\
n = 13:&\quad x^4 + 50x^3 + 123x^2 + 50x + 1,\\
n = 15:&\quad x^8 + 112x^7 - 630x^6 + 1568x^5 - 2109x^4\\
       &\qquad + 1568x^3 - 630x^2 + 112x + 1,\\
n = 25:&\quad x^8 + 13346x^7 + 24067x^6 - 4876x^5 + 11653x^4\\
       &\qquad - 4876x^3 + 24067x^2 + 13346x + 1 .
\end{align*}
\end{proposition}

The case $n = 7$ is emblematic: the phase of a level-$7$ Schmidt circle,
$\gamma$-normalized, \emph{is} the fundamental unit of $\Q(\sqrt3)$ up to
sign.  At the partially coherent levels the intermediate polynomials are
equally clean: for $w^3$, $x^2 + 6x + 1$ at $n = 5$ ($w^3 =
-\eps_8^{\pm2}$), $x^4 + 584x^3 + 766x^2 + 584x + 1$ at $n = 9$,
$x^4 + 25944604x^3 + 99499206x^2 + 25944604x + 1$ at $n = 17$; for $w^2$
at $n = 19$, $x^8 - 571772x^7 + 36149386x^6 - 42809072x^5 + 82773139x^4 -
(\text{sym})$.

The pattern of the table resists a congruence law: the sign part $m_2$
and the cube part $m_3$ of $m(n)$ vary independently: in the computed
range $m_2 = 2$ exactly at $n = 11$, $19$, $23$, $27$, $29$, $31$, $33$,
$35$, and $m_3 = 3$ exactly at $n = 5$, $9$, $11$, $17$, $21$, $23$, $27$,
$29$, $31$, $33$, $35$; and every candidate congruence law formulated on
the data through $n = 27$ was falsified at $29 \le n \le 35$.  The
mechanism is visible at the boundary forms, where the mirror law
genuinely fails by the explicit factors
\[
  (\gamma_2^2\gamma_3)(\tau + 1) = -\zeta_3\,(\gamma_2^2\gamma_3)(\tau),
  \qquad
  (\gamma_2^2\gamma_3)(-1/\tau) = -(\gamma_2^2\gamma_3)(\tau)
\]
(non-real $w$ on ambiguous boundary classes at $n = 31$); re-evaluating
the $\gamma$'s on a Schertz $6$-system, or using the $\eta$-quotient form
with a Hermite-adapted $\tau_0$, moves the failure elsewhere without
removing it.

\begin{openproblem}[the cocycle]\label{op:cocycle}
Determine $\zeta(\sigma, f)$ --- the multiplier bookkeeping of
Kubert--Lang and Schertz \cite[Ch.~11]{KubertLang1981}, \cite{Schertz2010}
for the hybrid object $w_f$, a class function $u_f$ times a ratio of level-$6$
class invariants --- and derive $m(n)$.  The certified table of
Proposition \ref{prop:mn} is the constraint any such computation must
reproduce; the $\eta^4$-form \eqref{eq:etaquot} suggests a Dedekind-sum
formula in $(r_0, s_0)$.
\end{openproblem}

%%% ============================================================
\section{The Robert index}\label{sec:robert}

\subsection{The quadratic layer}\label{subsec:quadratic}

\begin{proposition}[eigenprojections against fundamental units]
\label{prop:quadlayer}
Let $\chi$ be an odd real character of $\Cl(1 - n^2)$ with proved
factorization $L(s,\chi) = L(s,\chi_{d_1})L(s,\chi_{d_2})C(s)$
(Proposition \ref{prop:genus-hyp}), or a real character of $\Cl(\Oc_n)$
with the factorization of Proposition \ref{prop:genus-euc}, and set
\[
  m_\chi \;:=\; \frac{2h(d_1^*)}{w(d_1^*)}\,h(d_2)\,C(0) \in \Q .
\]
Then
\begin{gather*}
  \sum_f\chi(f)\log|R_f| = -24\,m_\chi\log\eps_{d_2}
  \qquad\text{(hyperbolic)},\\
  \sum_\idealc\chi(\idealc)\log|G_\idealc| = -12\,m_\chi\log\eps_{d_2}
  \qquad\text{(Euclidean)} :
\end{gather*}
the $\chi$-eigenprojection of the Schmidt-unit system against the
fundamental unit of its real quadratic field is an explicit integer
multiple of $\log\eps_{d_2}$ built from the class numbers $h(d_1^*)$,
$h(d_2)$ and the conductor multiplier $C(0)$.
\end{proposition}

\begin{proof}
Theorem \ref{thm:Rpoly}, resp.\ Theorem \ref{thm:Dn}(2), with
\eqref{eq:Lprime-genus}.
\end{proof}

The certified values (residuals $\le 10^{-248}$ at $250$ digits):
\begin{center}
\small
\begin{tabular}{@{}llll@{}}
\toprule
aspect & $n$ & $\chi$ & projection\\
\midrule
hyperbolic & 11 & $(0,1)$ & $-48\log\eps_8$\\
hyperbolic & 11 & $(1,0)$ & $-16\log\eps_{40}$\\
hyperbolic & 13 & $(0,1)$, $(1,0)$ & $-24\log\eps_{24}$, $-24\log\eps_{21}$\\
hyperbolic & 15 & $(0,1)$, $(2,1)$ & $-24\log\eps_{28}$, $-12\log\eps_{56}$\\
Euclidean & 3, 5, 7 & $\chi_2$ & $-4\log\eps_{12}$, $-24\log\eps_5$,
$-12\log\eps_{28}$\\
Euclidean & 9, 11, 13 & $\chi_2$ & $-16\log\eps_{12}$, $-12\log\eps_{44}$,
$-24\log\eps_{13}$\\
\bottomrule
\end{tabular}
\end{center}

\subsection{The cubic layer: the index is $8\,h_{L_3}$}
\label{subsec:cubic}

Among the Euclidean levels $n \le 31$, the ring class group $\Cl(\Oc_n)$
has $3$-torsion exactly at $n = 9, 11, 13, 18, 22, 23, 26, 27$.  At such a
level let $\chi_3$ be a cubic character, $F$ the cyclic cubic extension
of $K$ inside $H_n$ cut out by $\ker\chi_3$ --- a dihedral sextic over
$\Q$ --- and $L_3 \subset F$ its real cubic subfield, the fixed field of
$\ker\chi_3$ and the mirror.  Let
\[
  \theta \;:=\; \prod_{\idealc \in \ker\chi_3}G_\idealc
\]
be the principal $\Delta$-mass coset product --- at $n = 9, 11, 13$ the
real root of the coset cubics of Remark \ref{rem:cosetcubics}.  Since
$v_p(M(n))$ is a multiple of $3$ at every $p$ (the exponents of
\PaperI[Thm.~7.18] carry the factor $6/e_p \in \{3, 6\}$), $|M(n)|^{1/3}$
is an integer: $3^8$, $11^2$, $1$, $2^{12}3^{16}$, $2^{12}11^4$, $23^2$,
$2^{12}$, $3^{26}$ at the eight levels.

\begin{proposition}\label{prop:thetaunit}
The number $\theta_u := \theta/|M(n)|^{1/3}$ is an algebraic \emph{unit}
of $L_3$, and $\log|\theta_u| = -8\,L'(0,\chi_3)$ at the real embedding.
\end{proposition}

\begin{proof}
$\theta$ is fixed by the translations in $\ker\chi_3$ (Theorem
\ref{thm:Dn}(1)) and by complex conjugation (which permutes the coset
$\ker\chi_3$), so $\theta \in L_3$, real.  Integrality both ways: by
Theorem \ref{thm:perclass} the valuation of every $G_\idealc$ at every
place over $p \mid n$ is the constant $w_p(k)$ and $0$ elsewhere, so
$v(\theta) = \tfrac h3w_p(k) = \tfrac13v(M(n))$ over $p \mid n$ and $0$
elsewhere --- exactly the valuations of $|M(n)|^{1/3}$.  For the
logarithm write $a = \log|\theta|$ and $b$ for the common log-modulus of
the two conjugate coset products (they are complex conjugates of each
other, the mirror exchanging the two nontrivial cosets).  The uniform
Stark law $-12L'(0,\chi_3) = \sum\chi_3\log|G|$ of Theorem
\ref{thm:Dn}(2) reads $a - b = -12L'$ (as $\zeta_3 + \zeta_3^2 = -1$),
while $a + 2b = \log|M(n)|$; solving, $\log|\theta_u| = a -
\tfrac13\log|M(n)| = \tfrac23(a - b) = -8L'(0,\chi_3)$.
\end{proof}

\begin{lemma}[conductor Euler factors]\label{lem:conductor}
Let $\chi$ be a character of $\Cl(\Oc_n)$.  The Epstein $L$-function is
the sum $L(s,\chi) = \sum_\idealb\bar\chi(\idealb)N(\idealb)^{-s}$ over all
invertible integral $\Oc_n$-ideals, and factors as $L(s,\chi) =
\Lprim(s,\chi)\,C_n(s)$, where $\Lprim(s,\chi)$ is the Hecke $L$-function
of $\chi$ as a primitive ring class character of $K$ and $C_n(s) =
\prod_{p \mid n}C_{n,p}(s)$ collects the Euler factors of the
conductor-supported invertible ideals.  If $\chi$ is nontrivial on
$\ker(\Cl(\Oc_n) \to \Cl(\Oc_{n/p}))$ then $C_{n,p} \equiv 1$; in
particular $C_n \equiv 1$ when $\chi$ factors through no $\Cl(\Oc_m)$,
$m \mid n$, $m < n$.
\end{lemma}

\begin{proof}
For an invertible $\ideala$, $\alpha \mapsto (\alpha)\ideala^{-1}$
bijects the nonzero elements of $\ideala$ modulo $\pm1$ with the
invertible integral ideals in the class of $\ideala^{-1}$, and
$N(\alpha)/N\ideala = N((\alpha)\ideala^{-1})$; this gives the ideal sum,
and the Euler product over $p \nmid n$ is the usual one.  At $p \mid n$
the local order $\Oc_{n,p} = \Z_p + p^k\Oc_{K,p}$ is local with maximal
ideal $p\Z_p + p^k\Oc_{K,p}$, and its integral invertible ideals are the
$x\Oc_{n,p}$, $x \in \Oc_{n,p}\setminus\{0\}$ modulo $\Oc_{n,p}^\times$.
For a non-unit $x$ one checks $x\,\Oc_{n/p,p}^\times \subseteq \Oc_{n,p}$,
so the group $U' = \Oc_{n/p,p}^\times/\Oc_{n,p}^\times$ acts freely on
the non-unit ideals, preserving norms; the class of $u\Oc_{n,p}$ ($u \in
U'$, extended by $\Oc_{n,q}$ at $q \ne p$) runs through $\ker(\Cl(\Oc_n)
\to \Cl(\Oc_{n/p}))$, and $\chi$ is multiplicative along the orbit, so
every orbit sum carries the factor $\sum_{u \in U'}\bar\chi(u\Oc_n) = 0$.
Only the unit ideal survives: $C_{n,p} \equiv 1$.
\end{proof}

\begin{proposition}\label{prop:cubicfield}
$L_3$ has one real and one complex place (unit rank $1$), and
\[
  h_{L_3}\,R_{L_3}\,C_n(0) \;=\; L'(0,\chi_3),
\]
where $L(s,\chi_3)$ is the Epstein $L$-function of level $n$ and
$C_n(s)$ the finite Euler product of Lemma \ref{lem:conductor}, with
$C_n \equiv 1$ at $n = 9, 11, 13, 23$, where $\chi_3$ is primitive
(the class groups of the proper divisor levels have no $3$-torsion).
\end{proposition}

\begin{proof}
$\zeta_{L_3}(s) = \zeta(s)\Lprim(s,\chi_3)$ by the Artin formalism for
the dihedral extension $F/\Q$ ($\zeta_{L_3} = \zeta\cdot L(s,\rho)$ with
$\rho = \operatorname{Ind}_K^\Q\chi_3$ the two-dimensional
representation, and $L(s,\rho) = \Lprim(s,\chi_3)$), and $L(s,\chi_3) =
\Lprim(s,\chi_3)C_n(s)$ by Lemma \ref{lem:conductor}.  Since $L(0,\chi_3)
= 0$, $L'(0,\chi_3) \ne 0$ (Proposition \ref{prop:klf} and the
nonvanishing of $L(1,\chi_3)$) and $\zeta(0) = -\tfrac12$, the function
$\zeta_{L_3}$ vanishes to first order at $s = 0$: $r_1 + r_2 - 1 = 1$
with a real embedding present, forcing $(r_1, r_2) = (1, 1)$.  The class
number formula at $s = 0$ \cite[Cor.~VII.5.11]{Neukirch1999},
$\zeta'_{L_3}(0) = -h_LR_L/w_L$ with $w_L = 2$,
gives $h_LR_L = -2\zeta'_{L_3}(0) = \Lprim'(0,\chi_3)$, and $L'(0,\chi_3)
= C_n(0)\Lprim'(0,\chi_3)$ since $\Lprim(0,\chi_3) = 0$.
\end{proof}

\begin{lemma}[fundamental units; computational, rigorous]
\label{lem:fundunit}
For each of the eight levels, the real root $\eta_L$ of the polynomial
listed in Theorem \ref{thm:robert} is a fundamental unit of $L_3$,
$R_{L_3} = \log|\eta_L|$.  Method: starting from $\theta_u$, every $k$-th
root, $2 \le k \le 40$, is tested for certified integrality of its
minimal polynomial (absolute-error criterion, $\ge 249$ spare digits at
$250$), descending whenever a root lies in $L_3$; the descent terminates
at $\eta_L$.  If $\eta_L$ were a proper power $\eta_{\mathrm{true}}^{J}$,
$J > 1$, then some prime $k \mid J$ would satisfy $k \le J = R_{L_3}/
R_{\mathrm{true}} < R_{L_3}/0.2052 < 40$ by Friedman's unconditional
bound $R_K > 0.2052$ for every number field $K \ne \Q$
\cite{Friedman1989}, and $\eta_L^{1/k} \in L_3$ would have been detected.
The largest ratio is $R_{L_3}/0.2052 < 31.2$ (at $n = 23$).
\end{lemma}

\begin{theorem}[Robert index, cubic layer]\label{thm:robert}
At every Euclidean level with a cubic character, $n \in \{9, 11, 13, 18,
22, 23, 26, 27\}$,
\[
  \bigl[\Oc_{L_3}^\times : \langle -1,\ \theta_u\rangle\bigr]
  \;=\; \frac{\log|\theta_u|}{R_{L_3}}
  \;=\; 8\,\frac{L'(0,\chi_3)}{R_{L_3}}
  \;=\; 8\,h_{L_3}\,C_n(0),
\]
where $C_n(0) = 1$ when $\chi_3$ is primitive and $C_n(0)$ is the integer
imprimitivity multiplier of Lemma \ref{lem:conductor} otherwise.  The
identity is Propositions \ref{prop:thetaunit}--\ref{prop:cubicfield}; the
content of the theorem is the certified computation of the fundamental
units (Lemma \ref{lem:fundunit}), hence of $R_{L_3}$ and of the integer
$L'(0,\chi_3)/R_{L_3} = h_{L_3}C_n(0)$ ($\ge 249$ spare digits, $L'$ by
the independent evaluation of Lemma \ref{lem:gamma-eval}):
\begin{center}
\footnotesize
\setlength{\tabcolsep}{3.5pt}
\begin{tabular}{@{}rlllrrr@{}}
\toprule
$n$ & $\chi_3$ & fundamental unit of $L_3$ & $R_{L_3}$ & $h_{L_3}$ &
$C_n(0)$ & index\\
\midrule
9 & primitive & $x^3 + 15x^2 + 57x - 1$ & 4.0476417264 & 1 & 1 & \textbf{8}\\
11 & primitive & $x^3 - 25x^2 + 201x - 1$ & 5.3026856598 & 1 & 1 & \textbf{8}\\
13 & primitive & $x^3 - x^2 + 9x - 1$ & 2.1860813593 & 3 & 1 & \textbf{24}\\
23 & primitive & $x^3 - 49x^2 + 601x - 1$ & 6.3984592530 & 2 & 1 & \textbf{16}\\
18 & from $n = 9$ & $x^3 + 15x^2 + 57x - 1$ & 4.0476417264 & 1 & 2 & 16\\
22 & from $n = 11$ & $x^3 - 25x^2 + 201x - 1$ & 5.3026856598 & 1 & 2 & 16\\
26 & from $n = 13$ & $x^3 - x^2 + 9x - 1$ & 2.1860813593 & 3 & 2 & 48\\
27 & from $n = 9$ & $x^3 + 15x^2 + 57x - 1$ & 4.0476417264 & 1 & 4 & 32\\
\bottomrule
\end{tabular}
\end{center}
The minimal polynomials of $\theta_u$ (constant term $-1$; the
regulators are shown to $10$ decimals, and are certified to $250$):
{\footnotesize
\begin{align*}
n = 9:&\quad x^3 - 11708931x^2 + 115597311109635x - 1,\\
n = 11:&\quad x^3 + 2297078781x^2 + 2651044211389651971x - 1,\\
n = 13:&\quad x^3 + 446643445245x^2 + 61048319249786206560771x - 1,\\
n = 23:&\quad x^3 - 28994720086003708422147x^2\\
      &\qquad + 289100431655153833437856247634098764691761155x - 1,\\
n = 18:&\quad x^3 + 94095557056509x^2 + 13362738335777743394966415363x - 1,\\
n = 22:&\quad x^3 + 25517496658857981x^2
      + 7028035410742581725200408425098342403x - 1,\\
n = 26:&\quad x^3 - 77393728680750899988483x^2\\
      &\qquad + 3726897283223817102236830775918158257403004931x - 1,\\
n = 27:&\quad x^3 + 17871502813780746125713563645x^2\\
      &\qquad + 178562775830464135183026414005799966775800870888956534787x - 1 .
\end{align*}}
\end{theorem}

\begin{proof}
$\Oc_{L_3}^\times = \{\pm1\}\times\langle\eta_L\rangle$ by Proposition
\ref{prop:cubicfield} and Lemma \ref{lem:fundunit}, so the index of
$\langle -1, \theta_u\rangle$ is $\log|\theta_u|/\log|\eta_L|$, which is
$8L'(0,\chi_3)/R_{L_3}$ by Proposition \ref{prop:thetaunit} and
$8h_{L_3}C_n(0)$ by Proposition \ref{prop:cubicfield}.  At the primitive
levels $h_{L_3}$ is the certified integer $L'(0,\chi_3)/R_{L_3}$; at the
pullback levels the cubic field, hence $h_{L_3}$ and $R_{L_3}$, is that
of the primitive level, and the certified integer $L'(0,\chi_3)/R_{L_3}$
is $h_{L_3}C_n(0)$, from which $C_n(0)$ is read off.
\end{proof}

Three readings.  (i) At $n = 13$ and $n = 23$ the class numbers $h_{L_3}
= 3, 2$ of the cubic fields are forced into the index --- precisely the
Robert--Kubert--Lang phenomenon, ``cyclotomic units have index $h^+$'',
realized by the Schmidt disks of curvature $2n$.  The row $n = 23$ is a
genuine out-of-sample confirmation: a new cubic field, computed after
the law $8h_{L_3}$ had been formulated on $9, 11, 13$.  (ii) The pullback
levels are a strong internal consistency check: the \emph{identical}
fundamental units and regulators recur at $n = 9, 18, 27$, at $n = 11,
22$ and at $n = 13, 26$, while the coset units $\theta_u$ differ, and the
index grows by exactly the local Euler value $C_n(0)$ --- $2$ per added
prime, $4$ at the square-conductor step $9 \to 27$ --- the same local
Euler phenomenon as the conductor corrections of Section
\ref{sec:genus}.  (iii) The universal factor $8 = 12/(3/2)$ is the weight
bookkeeping of the $\Delta$-normalization: the coset object is a
$12$th-power-level datum, projected with the $\tfrac32$-coefficient of
the Stark relation of Remark \ref{rem:cosetcubics}.

\begin{conjecture}[Robert index for Schmidt units]\label{conj:robert}
For every Euclidean level whose cubic character is primitive,
$[\Oc_{L_3}^\times : \langle -1, \theta_u\rangle] = 8\,h_{L_3}$
(imprimitive levels: times the multiplier $C_n(0)$).  More generally, for
every ring class field $H_n$ of $\Q(i)$, the group generated by the
$\Delta$-data $\{G_\idealc\}$ modulo its mass normalization sits inside
the unit group of $H_n$ with index a product of class numbers of the
subfields cut out by the character group, up to powers of $2$ and $3$ ---
the shape of the Kubert--Lang index formulas
\cite[Ch.~12--13]{KubertLang1981}.
\end{conjecture}

The conjecture is certified at all eight computed cubic levels, and the
quadratic layer's class-number factors are Proposition
\ref{prop:quadlayer}.  We have not compared $\theta_u$ with the Stark
unit of $H_n$ \cite{Stark1980, HajirVillegas1997}, whose group has index
$h_{H_n}$ in $\Oc_{H_n}^\times$ (with the roots of unity) at prime
conductor; the two indices, $8h_{L_3}$ in $\Oc_{L_3}^\times$ and
$h_{H_n}$ in $\Oc_{H_n}^\times$, live in different groups, and relating
them is a natural next step outside the scope of this paper.

\begin{remark}[the hyperbolic side]\label{rem:hypindex}
The hyperbolic $R_f$ satisfy $R_{\idealr f} = 1/R_f$, so their
projections vanish on every character with $\chi(\idealr) = +1$; since
cubic characters are even on the $2$-torsion class $\idealr$, the
hyperbolic cubic layer degenerates identically.  The hyperbolic Robert
index lives entirely on the odd characters: the quadratic layer of
Proposition \ref{prop:quadlayer}, and the sextic $\chi_2\chi_3$-layer at
levels such as $n = 21$, where the odd-part coset products of the $R_f$
are units by Theorem \ref{thm:unit}; their index against the rank-$2$
unit lattice is the first genuinely hyperbolic index datum (Section
\ref{sec:outlook}).
\end{remark}

%%% ============================================================
\section{Outlook}\label{sec:outlook}

Four continuations are genuine, in the sense that the data needed to
start them is already computed and certified.

\begin{enumerate}
\item \emph{The cocycle of the first-power units} (Open problem
\ref{op:cocycle}).  Determine $\zeta(\sigma, f)$ by Shimura reciprocity
for $\gamma_2, \gamma_3$ (Gee--Stevenhagen, Schertz $N$-systems
\cite{Schertz2010}) applied to the hybrid object $w_f$; the certified
table $m(n)$ of Proposition \ref{prop:mn} is a sharp target, and the
$\eta^4$-form \eqref{eq:etaquot} suggests a Dedekind-sum formula in
$(r_0, s_0)$.
\item \emph{The genus-refined Gross--Zagier exponent law.}  The tables of
Propositions \ref{prop:dressing-euc} and \ref{prop:dressing-hyp} express
every genus-character $j$-sum as $k\log\eps_d + \sum_pe_p\lambda_p$ with
$e_p \in \{\pm3, \pm6, \pm9\}$ for $\Sigma_0$ and even for
$\Sigma_{1728}$, supported exactly on $\GZ(D,-3)$ resp.\ $\GZ(D,-4)$.
Gross--Zagier's $\epsilon(\mathfrak n)$-formula for the genus-character
weighted factorization of $(j - j')$-norms \cite{GrossZagier1985}, in
the conductor-degenerate refinement of Lauter--Viray
\cite{LauterViray2015}, should prove the exponents and signs.  This one
finite-place computation would settle at once the exact denominator
valuations of the phase \PaperI[\S6.7], the non-split part of the
$\lambda_n$-law (Corollary \ref{cor:inert} reduces it to the
$\beta$-collision slopes), and the general law of the conductor Euler
corrections at square conductors.
\item \emph{The hyperbolic sextic index.}  At $n = 21$, and at every
level whose class group has a $6$-part, the odd-character coset products
of the $R_f$ are units of the sextic subfield by Theorem \ref{thm:unit}
(Remark \ref{rem:hypindex}); their index against the rank-$2$ unit
lattice is the first genuinely hyperbolic Robert-index datum.
\item \emph{Stark recognition for characters of order $> 2$.}  The
certified coset cubics and quartics (Remark \ref{rem:cosetcubics},
Section \ref{subsec:nonfits}) pin $L'(0,\chi_3)$, $L'(0,\chi_4)$ as
logarithms of explicit algebraic numbers in the cubic and quartic
subfields of the ring class field; identifying them with the Stark units
of those fields --- a theorem over an imaginary quadratic base through
elliptic units \cite{Stark1980}, \cite[Ch.~6--7]{Schertz2010} --- would
turn the certified non-fits of Proposition \ref{prop:nonfits} into
positive closed forms, and is the natural route to Conjecture
\ref{conj:robert} beyond the cubic layer.
\end{enumerate}

\subsection*{What is new, with anchors}
Each item names the repository document and section where the full
record (proofs, data, script output) lives.
\begin{itemize}
\item The two Kronecker limit formulas for a geometric invariant of
circles: Theorems \ref{thm:euc-master} and \ref{thm:hyp-master}
(\texttt{phase-kronecker-\allowbreak limit.md}, \S2 and \S5).
\item The genus closed forms with the field $\Q(\sqrt n)$ and the exact
$\GZ$-supported dressings: Section \ref{sec:genus}
(\texttt{phase-kronecker-\allowbreak limit.md}, \S4 and \S6).
\item The unit theorem on all strata: Theorem \ref{thm:unit}
(\texttt{schmidt-\allowbreak units.md}, \S2).
\item The per-class valuation law and the split ladder: Theorem
\ref{thm:perclass} and Corollary \ref{cor:split}
(\texttt{schmidt-\allowbreak units.md}, \S3).
\item The first-power units, with $w = -\eps_{12}^{\pm1}$ at $n = 7$:
Section \ref{sec:firstpower} (\texttt{schmidt-\allowbreak units.md}, \S4).
\item The Robert index $8h_{L_3}C_n(0)$: Theorem \ref{thm:robert}
(\texttt{schmidt-\allowbreak units.md}, \S5).
\end{itemize}
Open, with the same anchors: the cocycle law $m(n)$
(\texttt{schmidt-\allowbreak units.md}, \S4); the exponent law of the dressings
(\texttt{phase-kronecker-\allowbreak limit.md}, \S8, open problems 1--2);
the sextic index (\texttt{schmidt-\allowbreak units.md}, \S5); and Stark recognition
(\texttt{phase-kronecker-\allowbreak limit.md}, \S8, open problem 3).

%%% ============================================================
\appendix

\section{Machine verification}\label{app:verification}

Every finite claim of this paper --- classwise laws, the limit-formula
identities against the independent evaluation, integer polynomial
coefficients, exact factorizations in quadratic fields, Newton polygons,
the coherence table, fundamental units and indices, and the certified
non-fits --- is verified by one of two standalone Python programs in the
repository accompanying the paper (\texttt{python3} with \texttt{mpmath}
\cite{mpmath} and \texttt{sympy} \cite{sympy}), each re-run from a clean
checkout for this draft; the statements imported from \PaperI{} (the
closed forms, $\omega \equiv 1$, the transport lemma, the $\Delta$-mass
law) are covered by the scripts catalogued in \PaperI[App.~A].  Runtimes
are on a single commodity core.

The certification guard rails of \PaperI[App.~A] are enforced
throughout: exact integer or rational arithmetic wherever a statement is
finite (Hermite normal forms, ideal arithmetic in quadratic fields with
Hensel-lifted valuations, exact factorization); a numerical value is
accepted as an integer or rational only with at least $\max(20,
\mathrm{dps}/5)$ spare digits in the absolute-error sense; working
precisions are set at run time, never at module import; and
integer-relation searches appear only in negative statements, in the
information-theoretically safe regime (a relation would be accepted only
if its information content were at most a quarter of the working
precision and its residual at rounding level).  No claim of the paper
rests on an integer-relation search.

\begin{center}
\footnotesize
\setlength{\tabcolsep}{3pt}
\renewcommand{\arraystretch}{0.94}
\begin{tabular}{@{}p{2.6cm}p{5.8cm}p{2.0cm}p{1.4cm}@{}}
\toprule
script / phase & verifies & precision, spare digits & time\\
\midrule
\texttt{phase\_klf.py --selftest} & Euclidean $n = 3, 5, 7, 9, 11, 13$
and hyperbolic $n = 9, 11, 13, 15$ at $250$ digits; Euclidean $n = 15$
at $400$ digits & & 4 min 39 s\\
\quad classwise laws & Lemmas \ref{lem:euc-classwise},
\ref{lem:hyp-classwise}, constants included, every class & dev.\
$\le 1.4\cdot10^{-248}$ ($1.4\cdot10^{-399}$ at 400d) & \\
\quad limit formula & $\sum\chi\log g = -12L'(0,\chi)$ against Lemma
\ref{lem:gamma-eval} ($600$--$4600$ terms), every character &
res.\ $\le 6.2\cdot10^{-249}$ & \\
\quad master identities & Theorems \ref{thm:euc-master},
\ref{thm:hyp-master}, every character; even hyperbolic sums vanish &
res.\ $\le 1.8\cdot10^{-248}$; $\le 7.7\cdot10^{-250}$ & \\
\quad anchor & \eqref{eq:euc-trivial} against certified $H_{-4n^2}$ and
the closed form of $M(n)$ & res.\ $\le 9\cdot10^{-249}$ & \\
\quad $D_n$ & Theorem \ref{thm:Dn}: integer coefficients, $D_n(0) =
(-1)^hM(n)$, uniform Stark law; coset cubics of Remark
\ref{rem:cosetcubics} & spare 226--248 (358 at $n = 15$) & \\
\quad genus factorizations & \eqref{eq:conv} as exact integer identities,
$m \le 300$; closed forms \eqref{eq:Lprime-genus} against Lemma
\ref{lem:gamma-eval}; hyperbolic table of Proposition
\ref{prop:genus-hyp}; even-character factorizations & exact; res.\
$\le 4.2\cdot10^{-250}$ & \\
\quad $\Sigma$-factorizations & coset products as integer quadratics;
every $A/B$ factorization re-multiplied exactly in $\Q(\sqrt d)$
(Propositions \ref{prop:dressing-euc}, \ref{prop:dressing-hyp}); GZ
support & spare 174--238 (315--317 at 400d) & \\
\quad non-fits & Proposition \ref{prop:nonfits}: safe PSLQ, height
$10^6$, tolerance $10^{-200}$ ($10^{-320}$ at 400d) & information
$\le\mathrm{dps}/4$ & \\
\midrule
\texttt{schmidt\_units.py --selftest} & four phases & & 4 min 0 s\\
\quad phase 1 & Lemma \ref{lem:lattice} classwise, odd $n \le 21$;
level polynomials of Proposition \ref{prop:Rrecord} (integer,
palindromic, constant term $1$, regression-checked, irreducible by
exact factorization); stratum polynomials and the $R \equiv 1$ collapse
of Proposition \ref{prop:strata} &
250--420d; dev.\ $\le 6.5\cdot10^{-419}$; spare 237--385 & \\
\quad phase 2 & Proposition \ref{prop:newton}: single-slope Newton
polygons of the certified $D_n$ at every $p \mid n$, $n \le 18, 21, 25,
27, 49$; split-ladder witness at $n = 15$; inert ladders of Corollary
\ref{cor:inert} re-derived from certified slopes & 250d (460d at $n =
49$); spare 140--249 & \\
\quad phase 3 & Proposition \ref{prop:wlaws}(1)--(3) at every class; the
table $m(n)$ of Proposition \ref{prop:mn}, each entry re-certified;
coherent-level polynomials & 250--500d & \\
\quad phase 4 & Proposition \ref{prop:quadlayer} projections (res.\
$\le 6.2\cdot10^{-249}$); at all eight cubic levels: unit-normalized coset
cubics, root descent $k \le 40$ with the Friedman cap, $L'(0,\chi_3)/
R_{L_3}$ as certified integers, $\log|\theta_u| = -8L'$, the index, and
the regression record of Theorem \ref{thm:robert} & 250d; spare
249--250 & 14--43 s per level\\
\bottomrule
\end{tabular}
\end{center}

%% TODO before submission: repository URL / archival identifier for
%% the scripts.

%%% ============================================================
\begin{thebibliography}{99}

\bibitem{Cox2013}
D.~A.~Cox,
\emph{Primes of the form $x^2 + ny^2$: Fermat, class field theory,
and complex multiplication},
2nd ed., Wiley, 2013.

\bibitem{DiamondShurman2005}
F.~Diamond, J.~Shurman,
\emph{A first course in modular forms},
Graduate Texts in Mathematics 228, Springer, 2005.

\bibitem{Friedman1989}
E.~Friedman,
\emph{Analytic formulas for the regulator of a number field},
Invent.\ Math.\ \textbf{98} (1989), 599--622.

\bibitem{GillardRobert1979}
R.~Gillard, G.~Robert,
\emph{Groupes d'unit\'es elliptiques},
Bull.\ Soc.\ Math.\ France \textbf{107} (1979), 305--317.

\bibitem{Gross1986}
B.~H.~Gross,
\emph{On canonical and quasi-canonical liftings},
Invent.\ Math.\ \textbf{84} (1986), 321--326.

\bibitem{GrossZagier1985}
B.~H.~Gross, D.~B.~Zagier,
\emph{On singular moduli},
J.\ reine angew.\ Math.\ \textbf{355} (1985), 191--220.

\bibitem{HajirVillegas1997}
F.~Hajir, F.~Rodriguez Villegas,
\emph{Explicit elliptic units, I},
Duke Math.\ J.\ \textbf{90} (1997), 495--521.

\bibitem{SchmidtI}
%% TODO: update on publication of Paper I.
K.~Ivar,
\emph{Schmidt circles of the Gaussian integers: classification, class
numbers, and the phase invariant},
preprint (2026); Paper I of this series.

\bibitem{mpmath}
F.~Johansson et al.,
\emph{mpmath: a Python library for arbitrary-precision floating-point
arithmetic}, \texttt{https://mpmath.org}.

\bibitem{KubertLang1981}
D.~S.~Kubert, S.~Lang,
\emph{Modular units},
Grundlehren der mathematischen Wissenschaften 244, Springer, 1981.

\bibitem{Kucuksakalli2012}
\"O.~K\"u\c{c}\"uksakall{\i},
\emph{Class numbers of ring class fields of prime conductor},
Acta Arith.\ \textbf{153} (2012), 251--269.

\bibitem{Lang1987}
S.~Lang,
\emph{Elliptic functions},
2nd ed., Graduate Texts in Mathematics 112, Springer, 1987.

\bibitem{LauterViray2015}
K.~Lauter, B.~Viray,
\emph{On singular moduli for arbitrary discriminants},
Int.\ Math.\ Res.\ Not.\ IMRN \textbf{2015}, no.\ 19, 9206--9250.

\bibitem{sympy}
A.~Meurer et al.,
\emph{SymPy: symbolic computing in Python},
PeerJ Comput.\ Sci.\ \textbf{3}:e103 (2017).

\bibitem{Meyer1957}
C.~Meyer,
\emph{Die Berechnung der Klassenzahl abelscher K\"orper \"uber
quadratischen Zahlk\"orpern},
Akademie-Verlag, Berlin, 1957.

\bibitem{Neukirch1999}
J.~Neukirch,
\emph{Algebraic number theory},
Grundlehren der mathematischen Wissenschaften 322, Springer, 1999.

\bibitem{Oukhaba2003}
H.~Oukhaba,
\emph{Index formulas for ramified elliptic units},
Compositio Math.\ \textbf{137} (2003), 1--22.

\bibitem{Ramachandra1964}
K.~Ramachandra,
\emph{Some applications of Kronecker's limit formulas},
Ann.\ of Math.\ (2) \textbf{80} (1964), 104--148.

\bibitem{Robert1973}
G.~Robert,
\emph{Unit\'es elliptiques},
Bull.\ Soc.\ Math.\ France, M\'em.\ \textbf{36} (1973).

\bibitem{Schertz2010}
R.~Schertz,
\emph{Complex multiplication},
New Mathematical Monographs 15, Cambridge University Press, 2010.

\bibitem{Schmidt1975}
A.~L.~Schmidt,
\emph{Diophantine approximation of complex numbers},
Acta Math.\ \textbf{134} (1975), 1--85.

\bibitem{Schoeneberg1939}
B.~Schoeneberg,
\emph{Das Verhalten von mehrfachen Thetareihen bei
Modulsubstitutionen},
Math.\ Ann.\ \textbf{116} (1939), 511--523.

\bibitem{Siegel1961}
C.~L.~Siegel,
\emph{Lectures on advanced analytic number theory},
notes by S.~Raghavan, Tata Institute of Fundamental Research, Bombay,
1961.

\bibitem{StangeTAMS}
K.~E.~Stange,
\emph{The Apollonian structure of Bianchi groups},
Trans.\ Amer.\ Math.\ Soc.\ \textbf{370} (2018), 6169--6219.

\bibitem{StangeIMRN}
K.~E.~Stange,
\emph{Visualising the arithmetic of imaginary quadratic fields},
Int.\ Math.\ Res.\ Not.\ IMRN \textbf{2018}, no.\ 12, 3908--3938.

\bibitem{Stark1980}
H.~M.~Stark,
\emph{$L$-functions at $s = 1$. IV. First derivatives at $s = 0$},
Adv.\ in Math.\ \textbf{35} (1980), 197--235.

\bibitem{Sturm1987}
J.~Sturm,
\emph{On the congruence of modular forms},
in: Number theory (New York, 1984--1985), Lecture Notes in Math.\ 1240,
Springer, 1987, 275--280.

\bibitem{VlasenkoZagier2013}
M.~Vlasenko, D.~Zagier,
\emph{Higher Kronecker ``limit'' formulas for real quadratic fields},
J.\ reine angew.\ Math.\ \textbf{679} (2013), 23--64.

\end{thebibliography}

\end{document}
